\documentclass[class=article, crop=false, oneside]{standalone}
% ==============================================================================
% SETUP.TEX - CONFIGURACIÓN GLOBAL DEL PROYECTO TC2
% ==============================================================================
\ifdefined\SETUPCONFIGURED
\else
\newcommand{\SETUPCONFIGURED}{}

% Este archivo centraliza todos los paquetes y estilos.
% Debe incluirse en el preámbulo del libro principal y de los archivos standalone.

% --- 1. CODIFICACIÓN E IDIOMA ---
\usepackage[utf8]{inputenc}
\usepackage[T1]{fontenc}
\usepackage[spanish, es-tabla]{babel}  
\usepackage{csquotes} 

\usepackage{import}
\usepackage[mode=image, subpreambles=false]{standalone}

% --- 2. MATEMÁTICAS Y CIENCIAS ---
\usepackage{amsmath, amssymb, amsfonts, mathtools}

% Operadores en notación español/europea (seno, coseno y seno hiperbólicos)
\DeclareMathOperator{\sen}{sen}
\DeclareMathOperator{\ch}{ch}
\DeclareMathOperator{\sh}{sh}

% --- 3. DISEÑO Y GEOMETRÍA --- 
\usepackage[left=2.5cm,right=2.5cm,top=3cm,bottom=3cm]{geometry}
\usepackage{parskip} 
\usepackage{float}   

% --- 4. GRÁFICOS Y TABLAS ---
\usepackage{graphicx}
\usepackage{booktabs} 
\usepackage{tikz}
\usetikzlibrary{arrows.meta, calc, angles, quotes, babel, shapes.geometric}
\usepackage{pgfplots}
\usepgfplotslibrary{groupplots}
\usepgfplotslibrary{external}
\usepgfplotslibrary{patchplots}
\pgfplotsset{compat=1.18}
\usepackage[siunitx, RPvoltages]{circuitikz}

% --- 5. COLORES Y CAJAS ---
\usepackage[table]{xcolor}
\usepackage[theorems, most]{tcolorbox}

\definecolor{utnblue}{RGB}{0, 51, 102}
\definecolor{codegreen}{rgb}{0,0.6,0}
\definecolor{codegray}{rgb}{0.5,0.5,0.5}
\definecolor{codepurple}{rgb}{0.58,0,0.82}
\definecolor{backcolour}{rgb}{0.96,0.96,0.96}

% --- 6. ENTORNOS PERSONALIZADOS (CAJAS) ---
\newtcolorbox{pedagogico}{
    colback=codegreen!5!white,
    colframe=codegreen,
    title=\textbf{Enfoque Pedagógico},
    fonttitle=\bfseries,
    sharp corners,
    boxrule=0.5mm
}

\newtcolorbox{ingenieria}{
    colback=utnblue!5!white,
    colframe=utnblue,
    title=\textbf{Utilidad Ingenieril},
    fonttitle=\bfseries,
    sharp corners,
    boxrule=0.5mm
}

\newtcolorbox{nota}{
    colback=codegray!5!white,
    colframe=codegray,
    title=\textbf{Nota},
    fonttitle=\bfseries,
    sharp corners,
    boxrule=0.5mm
}

% --- 7. CÓDIGO FUENTE (LISTINGS) ---
\usepackage{listings}
\lstdefinestyle{miestilo}{
    backgroundcolor=\color{backcolour},   
    commentstyle=\color{codegreen},
    keywordstyle=\color{magenta},
    numberstyle=\tiny\color{codegray},
    stringstyle=\color{codepurple},
    basicstyle=\ttfamily\footnotesize,
    breakatwhitespace=false,         
    breaklines=true,                 
    captionpos=b,                    
    keepspaces=true,                 
    numbers=left,                    
    numbersep=5pt,                  
    showspaces=false,                
    showstringspaces=false,
    showtabs=false,                  
    tabsize=2,
    frame=single,                    
    rulecolor=\color{codegray}
}
\lstset{style=miestilo}
\lstset{literate={ñ}{{\~n}}1 {á}{{\'a}}1 {é}{{\'e}}1 {í}{{\'i}}1 {ó}{{\'o}}1 {ú}{{\'u}}1}

% --- 8. HIPERVÍNCULOS ---
\usepackage[colorlinks=true, linkcolor=utnblue, urlcolor=utnblue, citecolor=utnblue]{hyperref}

% --- 9. PLACEHOLDER PARA FIGURAS PENDIENTES ---
% Marca el lugar y la descripción de una figura aún no dibujada (p. ej. Cap. 5:
% circuitos y gráficas diferidos). Greppable por \figpendiente para el pase posterior.
\newcommand{\figpendiente}[1]{%
  \begin{center}%
  \setlength{\fboxsep}{6pt}%
  \fbox{\parbox{0.85\linewidth}{\centering\small\itshape
  [\textbf{Figura pendiente}]\\[2pt] #1}}%
  \end{center}}

\fi

% ------------------------------------------------------------------
% Trabajo Práctico Nº 4 — Filtros: Teoría Imagen
% PARTE II: Resolución completa. Cada ejercicio incluye su enunciado
% completo (caja  "Enunciado") seguido del desarrollo.
% (El enunciado y las respuestas de la guía están en
%  TC-TP4-24-enunciado.tex.)
% Basada en la teoría del Capítulo 4 (Filtros — Teoría Imagen).
% Cátedra Teoría de Circuitos 2 — UTN-FRM — C. J. Larriqueta.
% ------------------------------------------------------------------
 
% Encabezado/pie de página
\usepackage{fancyhdr}
\setlength{\headheight}{14pt}
\pagestyle{fancy}
\fancyhf{}
\lhead{\small\color{utnblue} Teoría de Circuitos 2 --- TP Nº 4}
\rhead{\small\color{utnblue} Mg. Ing. Javier Velez}
\cfoot{\small\thepage}
\renewcommand{\headrulewidth}{0.4pt}

% Macro de ejercicio (encabezado uniforme)
\newcounter{ejer}
\newcommand{\ejercicio}[1]{%
  \refstepcounter{ejer}%
  \vspace{1.5ex}%
  \noindent{\large\bfseries\color{utnblue} Ejercicio \theejer.}\enspace{\bfseries #1}\par\nopagebreak\vspace{0.5ex}}

% Bloque (sin recuadro) para mostrar el enunciado completo de cada ejercicio
\newenvironment{enunciadobox}{%
    \par\addvspace{0.6ex}\begingroup\small\leftskip=1.2em\rightskip=1.2em
    \noindent{\bfseries\color{utnblue}Enunciado.}\enspace\ignorespaces}{%
    \par\endgroup\addvspace{0.8ex}}

\begin{document}

% ===================== PORTADA / TÍTULO =====================
\begin{center}
    {\Large\bfseries Universidad Tecnológica Nacional --- F.R. Mendoza}\\[2pt]
    {\large Teoría de Circuitos 2}\\[10pt]
    {\huge\bfseries\color{utnblue} Resolución de Trabajo Práctico Nº 4}\\[4pt]
    {\LARGE\bfseries Filtros: Teoría Imagen}\\[10pt]
\end{center}

\vspace{0.5ex}
\hrule
\vspace{1ex}

\begin{nota}
La resolución que sigue se desarrolla íntegramente con la teoría del
\textbf{Capítulo 4 --- Filtros (Teoría Imagen)}: condición de corte
$X_1=-4X_2$, impedancias características $z_{0T}$ y $z_{0\pi}$, diseño del
prototipo ($z_1 z_2 = k^2$), $m$-derivación serie, hemisección de adaptación
($m=0{,}6$), normalización/desnormalización y transformaciones de frecuencia.
Los valores obtenidos se contrastan con las respuestas oficiales de la guía.
\end{nota}

\subsection*{Marco teórico (formulario del Capítulo 4)}

Toda la resolución se apoya en la sección no disipativa (cuadripolo $L$-$C$ en
configuración T) del Capítulo 4. Recordamos los resultados que usaremos:

\begin{itemize}\setlength{\itemsep}{2pt}
    \item \textbf{Condición de corte:} $X_1 = -4X_2$, siendo $X_1$ la reactancia
    de la rama serie y $X_2$ la de la rama derivación.
    \item \textbf{Impedancias características:}
    \[
        z_{0T} = \sqrt{z_1 z_2 + \frac{z_1^2}{4}}\ ,\qquad
        z_{0\pi} = \frac{z_1 z_2}{z_{0T}}\ .
    \]
    \item \textbf{Prototipo pasabajos, sección T} (producto $z_1 z_2 = k^2$ constante)\textbf{:}
    \[
        \omega_c = \frac{2}{\sqrt{L_1 C_2}}\ ,\qquad
        R = \sqrt{\frac{L_1}{C_2}} = R_L\ \Longrightarrow\
        \boxed{\,L_1 = \frac{2R}{\omega_c}\,,\quad C_2 = \frac{2}{R\,\omega_c}\,}
    \]
    La rama serie se reparte en dos brazos de $L_1/2$ cada uno.
    \item \textbf{$m$-derivación serie:} $X_1' = m X_1$,\;
    $X_2' = \dfrac{X_2}{m} + X_1\dfrac{1-m^2}{4m}$, con $0<m<1$. La rama
    derivación es un resonante serie $\left[L_1\dfrac{1-m^2}{4m}\right]$ con
    $\left[m C_2\right]$, y
    \[
        \omega_\infty = \frac{\omega_c}{\sqrt{1-m^2}}\ \Longleftrightarrow\
        \boxed{\,m = \sqrt{1-\left(\frac{\omega_c}{\omega_\infty}\right)^2}\,}
    \]
    \item \textbf{Hemisección de adaptación} ($m = 0{,}6$): rama serie
    $m\,L_1/2$; rama derivación = $\left[L_1\dfrac{1-m^2}{2m}\right]$ en serie con
    $\left[C_2\dfrac{m}{2}\right]$. Su $z_{0\pi}'$ es casi constante e igual a
    $\sqrt{L_1/C_2}$ en el 85\,\% de la banda.
    \item \textbf{Normalización:} se fijan dos normas \emph{adimensionales}, de
    impedancia $a$ y de frecuencia $N$, dividiendo cada referencia por su unidad:
    \[
        a = \frac{R_L}{1\,\Omega}\ ,\qquad
        N = \frac{\omega_c}{1\,\tfrac{\text{rad}}{\text{s}}}\ \text{(PB/PA)}\quad\text{ó}\quad
        N = \frac{\omega_2-\omega_1}{1\,\tfrac{\text{rad}}{\text{s}}}\ \text{(PBanda/SBanda)} .
    \]
    Al ser las normas adimensionales, los elementos \emph{normalizados conservan
    sus unidades} (henrios y faradios). Se normaliza y se desnormaliza con:
    \[
        \bar L = L\,\frac{N}{a}\ ,\quad \bar C = C\,a\,N\ ,\quad \bar R = \frac{R}{a}
        \qquad\Longleftrightarrow\qquad
        L = \frac{\bar L\,a}{N}\ ,\quad C = \frac{\bar C}{a\,N}\ ,\quad R = \bar R\,a\ .
    \]
    \item \textbf{Transformación a pasaaltos (inversión $s = 1/s'$):} invierte el
    eje de frecuencias, de modo que la banda pasante del pasabajos
    ($0<\omega'<\omega_c$) se mapea en la del pasaaltos ($\omega_c<\omega<\infty$).
    En los elementos: un inductor $s'L\to \dfrac{L}{s}=\dfrac{1}{s\,(1/L)}$, esto es
    un \emph{capacitor} $1/L$; un capacitor $\dfrac{1}{s'C}\to \dfrac{s}{C}=s\,(1/C)$,
    esto es un \emph{inductor} $1/C$. En síntesis, cada inductor se vuelve capacitor
    y cada capacitor inductor, con valor recíproco.
    \item \textbf{Transformación cuadrática (pasabanda desde PB, suprimebanda desde
    PA):} $s' = s + \dfrac{\omega_0^2}{s}$, con $\omega_0 = \sqrt{\omega_1\omega_2}$
    el centro \emph{geométrico} de la banda y ancho de banda normalizado $=1$ (por
    eso $\omega_1\omega_2=\omega_0^2$). El centro de la banda del prototipo
    ($\omega'=0$) se mapea en $\omega=\omega_0$ y los cortes $\omega'=\pm1$ en
    $\omega_1$ y $\omega_2$. La clave es \emph{dónde} se aplica la sustitución: la
    rama serie se describe por su \underline{impedancia} y la derivación por su
    \underline{admitancia}, y en ambos casos $s'$ se reemplaza en una suma de dos
    términos, lo que genera una rama resonante:
    \begin{itemize}\setlength{\itemsep}{1pt}
        \item rama serie ---inductor---: $\;z = s'L \to sL + \dfrac{\omega_0^2 L}{s}
        = sL + \dfrac{1}{s\,\left(1/\omega_0^2 L\right)}$, o sea $L$ en \emph{serie}
        con $C=\dfrac{1}{\omega_0^2 L}$ $\Rightarrow$ resonante \emph{serie};
        \item rama derivación ---capacitor---: $\;y = s'C \to sC + \dfrac{\omega_0^2 C}{s}
        = sC + \dfrac{1}{s\,\left(1/\omega_0^2 C\right)}$, o sea $C$ en \emph{paralelo}
        con $L=\dfrac{1}{\omega_0^2 C}$ $\Rightarrow$ resonante \emph{paralelo}.
    \end{itemize}
    Cada par $L$-$C$ resuena en $\omega_0$ (pues $1/\sqrt{LC}=\omega_0$),
    reproduciendo en torno a $\omega_0$ el comportamiento que el elemento original
    tenía en el centro de la banda del prototipo. La regla (inductor$\to$serie,
    capacitor$\to$paralelo) es \emph{universal}: para el suprimebanda se aplica la
    \emph{misma} transformación al pasaaltos, y como éste tiene $C$ en serie y $L$
    en derivación, la rama serie resulta un resonante \emph{paralelo} y la
    derivación uno \emph{serie} (no cambia la regla, cambia el prototipo de partida).
\end{itemize}

\begin{nota}
\textit{Convención de la rama serie.} El libro reporta directamente el valor
\emph{colocado} en cada brazo serie ($L_1/2$ en el prototipo, $m\,L_1/2$ en la
$m$-derivada y en la hemisección). Las respuestas oficiales de la guía, en cambio,
tabulan la rama serie \underline{sin} dividir (la reactancia completa $L_1$ ó
$m\,L_1$) y aclaran que ``los valores finales de la rama serie se obtienen
dividiendo por dos''. Conviene tener presente la diferencia: la rama serie de una
\textbf{sección completa} (T) es $z_1$, \emph{repartida} en dos brazos de $z_1/2$
(la guía tabula $z_1$, el valor sin dividir; cada brazo colocado es $z_1/2$). Esos
dos brazos \emph{no} están en serie entre sí: entre ellos se conecta la rama
derivación, que deriva parte de la corriente. En una \textbf{media sección}
(hemisección) hay un \emph{único} brazo $z_1/2$, y el valor tabulado por la guía es
sólo el doble del colocado.
\end{nota}

% --------------------------------------------------------------
\ejercicio{Esquema de un filtro completo (Teoría Imagen).}

\begin{enunciadobox}
Esquematizar el diseño de un filtro completo según la Teoría Imagen.
\end{enunciadobox}

Un filtro completo se arma en cascada con \emph{adaptación imagen} entre
secciones. El núcleo es \textbf{al menos una sección prototipo} (da la atenuación
lejos del corte); se agrega una \textbf{sección $m$-derivada} (con $m$ pequeño)
para estrechar la transición cerca de $\omega_c$, y \textbf{no pueden faltar las
dos hemisecciones de adaptación} ($m = 0{,}6$). Cada hemisección es media
sección $L$, un cuadripolo \emph{asimétrico} con una impedancia imagen distinta
en cada puerto: por su \textbf{cara T} (impedancia imagen $z_{0T}$) empalma con
la cadena interna de secciones --- prototipo y $m$-derivada van en configuración
T y presentan $z_{0T}$ en sus bornes ---, y por su \textbf{cara $\pi$}
(impedancia imagen $z_{0\pi}$) se conecta a las resistencias de terminación
$R_L$. El valor $m = 0{,}6$ es precisamente el que vuelve $z_{0\pi}$ casi
constante e igual a $R=\sqrt{L_1/C_2}$ en el grueso de la banda (mientras que
$z_{0T}$ del prototipo varía mucho con la frecuencia); así la cara $\pi$ adapta
bien a la resistencia real y la cara T mantiene la continuidad de impedancia
imagen hacia adentro:

\begin{figure}[H]
    \centering
    \resizebox{\textwidth}{!}{% Figura generada con netlist2tikz (n2t render fig_ej01_bloques.sch --tikz --no-standalone).
% Netlist fuente (fig_ej01_bloques.sch):
%   # Esquema de filtro completo: Eg + RG -- [4 cuadripolos en cascada] -- RL
%   # up/down=1.30 centra las cajas (sobresalen igual arriba y abajo).
%   V g0 v0; up=1.30, l=E_g
%   W v0 v1; right=0.3
%   RG v1 a_in; right=1.0, l=R_G
%   W g0 a_inb; right=1.3
%   TP1 b1t b1b a_in a_inb; right=1.3, fill=utnblue!8, l={\shortstack{Hemisec.\\adaptación\\$m{=}0{,}6$}}
%   TP2 b2t b2b b1t b1b; right=1.3, fill=utnblue!8, l=Prototipo
%   TP3 b3t b3b b2t b2b; right=1.3, fill=utnblue!8, l={\shortstack{Sección\\$m$-derivada}}
%   TP4 b4t b4b b3t b3b; right=1.3, fill=utnblue!8, l={\shortstack{Hemisec.\\adaptación\\$m{=}0{,}6$}}
%   W b4t r1; right=0.4
%   W b4b r1b; right=0.4
%   RL r1 r1b; down=1.30, l=R_L
%
\begin{tikzpicture}[scale=1.00, transform shape, /tikz/circuitikz/bipoles/length=1.50cm, american currents, american voltages, font={\fontsize{11}{9}\selectfont}, voltage dir=RP]
  \coordinate (g0) at (0,0);
  \coordinate (v0) at (0,1.95);
  \coordinate (v1) at (0.45,1.95);
  \coordinate (a_in) at (1.95,1.95);
  \coordinate (a_inb) at (1.95,0);
  \coordinate (TP1@wnw) at (2.6,1.95);
  \coordinate (TP1@wsw) at (2.6,0);
  \coordinate (TP1@ene) at (5.2,1.95);
  \coordinate (TP1@ese) at (5.2,0);
  \coordinate (TP1@mid) at (3.9,0.975);
  \coordinate (b1t) at (5.85,1.95);
  \coordinate (b1b) at (5.85,0);
  \coordinate (TP2@wnw) at (6.5,1.95);
  \coordinate (TP2@wsw) at (6.5,0);
  \coordinate (TP2@ene) at (9.1,1.95);
  \coordinate (TP2@ese) at (9.1,0);
  \coordinate (TP2@mid) at (7.8,0.975);
  \coordinate (b2t) at (9.75,1.95);
  \coordinate (b2b) at (9.75,0);
  \coordinate (TP3@wnw) at (10.4,1.95);
  \coordinate (TP3@wsw) at (10.4,0);
  \coordinate (TP3@ene) at (13,1.95);
  \coordinate (TP3@ese) at (13,0);
  \coordinate (TP3@mid) at (11.7,0.975);
  \coordinate (b3t) at (13.65,1.95);
  \coordinate (b3b) at (13.65,0);
  \coordinate (TP4@wnw) at (14.3,1.95);
  \coordinate (TP4@wsw) at (14.3,0);
  \coordinate (TP4@ene) at (16.9,1.95);
  \coordinate (TP4@ese) at (16.9,0);
  \coordinate (TP4@mid) at (15.6,0.975);
  \coordinate (b4t) at (17.55,1.95);
  \coordinate (b4b) at (17.55,0);
  \coordinate (r1) at (18.15,1.95);
  \coordinate (r1b) at (18.15,0);
  \draw (v0) to [V, l_=$E_g$, n=V] (g0);
  \draw (v0) to (v1);
  \draw (v1) to [R, l_=$R_G$, n=RG] (a_in);
  \draw (g0) to (a_inb);
  \draw (3.9,0.975) node[rectangle, thick, inner sep=0pt, minimum width=2.60cm, minimum height=2.60cm, text width=2.08cm, align=center, shape border rotate=0, fill=utnblue!8, draw, scale=1.0] (TP1) {\shortstack{Hemisec.\\adaptación\\$m{=}0{,}6$}};
  \draw (a_in) -- (TP1@wnw);
  \draw (a_inb) -- (TP1@wsw);
  \draw (TP1@ene) -- (b1t);
  \draw (TP1@ese) -- (b1b);
  \draw (7.8,0.975) node[rectangle, thick, inner sep=0pt, minimum width=2.60cm, minimum height=2.60cm, text width=2.08cm, align=center, shape border rotate=0, fill=utnblue!8, draw, scale=1.0] (TP2) {Prototipo};
  \draw (b1t) -- (TP2@wnw);
  \draw (b1b) -- (TP2@wsw);
  \draw (TP2@ene) -- (b2t);
  \draw (TP2@ese) -- (b2b);
  \draw (11.7,0.975) node[rectangle, thick, inner sep=0pt, minimum width=2.60cm, minimum height=2.60cm, text width=2.08cm, align=center, shape border rotate=0, fill=utnblue!8, draw, scale=1.0] (TP3) {\shortstack{Sección\\$m$-derivada}};
  \draw (b2t) -- (TP3@wnw);
  \draw (b2b) -- (TP3@wsw);
  \draw (TP3@ene) -- (b3t);
  \draw (TP3@ese) -- (b3b);
  \draw (15.6,0.975) node[rectangle, thick, inner sep=0pt, minimum width=2.60cm, minimum height=2.60cm, text width=2.08cm, align=center, shape border rotate=0, fill=utnblue!8, draw, scale=1.0] (TP4) {\shortstack{Hemisec.\\adaptación\\$m{=}0{,}6$}};
  \draw (b3t) -- (TP4@wnw);
  \draw (b3b) -- (TP4@wsw);
  \draw (TP4@ene) -- (b4t);
  \draw (TP4@ese) -- (b4b);
  \draw (b4t) to (r1);
  \draw (b4b) to (r1b);
  \draw (r1) to [R, l_=$R_L$, n=RL] (r1b);
\end{tikzpicture}
}
\end{figure}

\noindent Cada bloque es un cuadripolo $L$-$C$ no disipativo. Al adaptar por
impedancia imagen no hay reflexiones en las uniones, de modo que las
\emph{funciones de propagación se suman}: con $\theta=\alpha+j\beta$,
$\theta_{\text{total}}=\theta_1+\dots+\theta_n$ y, por lo tanto, las atenuaciones
$\alpha$ (en dB/Np) se suman. Esto es la transferencia imagen en cascada del
Capítulo~3,
\[
    G(s)=\frac{V_2}{V_1}\frac{V_3}{V_2}\dotsm\frac{V_{n+1}}{V_n}
    =\sqrt{\frac{z_{In+1}}{z_{I1}}}\;e^{-(\theta_1+\theta_2+\dotsb+\theta_n)} :
\]
las relaciones de tensión se \emph{multiplican} (y las raíces de impedancia se
cancelan telescópicamente), lo que en el exponente equivale a \emph{sumar} los
$\theta_i$. Así se combina la buena transición de la $m$-derivada con la buena
atenuación lejana del prototipo.

% --------------------------------------------------------------
\ejercicio{Prototipo pasabajos T: $\omega_c = 20.000\ \tfrac{\text{rad}}{\text{s}}$, $R_L = 500\ \Omega$.}

\begin{enunciadobox}
Calcular el prototipo de un filtro pasabajo de configuración ``T'' con los
siguientes parámetros: $\omega_c = 20.000\ \tfrac{\text{rad}}{\text{s}}$ y
$R_L = 500\ \Omega$.
\end{enunciadobox}

Lo resolvemos \emph{desde las ecuaciones originales} del Capítulo~4, sin partir
de la fórmula ya despejada. En la sección prototipo pasabajos T las reactancias
de la rama serie y la rama derivación son
\[
    X_1 = \omega L_1\ ,\qquad X_2 = -\frac{1}{\omega C_2}\ .
\]

\noindent\textbf{(i) Condición de corte} ($X_1 = -4X_2$). La frecuencia de corte
es aquella para la cual $X_1=-4X_2$:
\[
    \omega_c L_1 = \frac{4}{\omega_c C_2}
    \ \Longrightarrow\
    \omega_c^2 = \frac{4}{L_1 C_2}
    \ \Longrightarrow\
    \boxed{\,L_1 C_2 = \frac{4}{\omega_c^2}\,}
    \quad\left(\text{equivale a }\omega_c = \frac{2}{\sqrt{L_1 C_2}}\right).
\]

\noindent\textbf{(ii) Impedancia característica en la banda} (adaptación a $R_L$).
La impedancia característica de la sección T,
\[
    z_{0T} = \sqrt{z_1 z_2 + \frac{z_1^2}{4}}
           = \sqrt{\frac{L_1}{C_2} - \frac{\omega^2 L_1^2}{4}}
           = \sqrt{\frac{L_1}{C_2}}\,\sqrt{1-\left(\frac{\omega}{\omega_c}\right)^2}\ ,
\]
vale $\sqrt{L_1/C_2}$ en el extremo de banda ($\omega\to 0$). La teoría imagen
exige adaptar esa impedancia nominal a la resistencia de terminación, de modo que
\[
    \boxed{\,\sqrt{\frac{L_1}{C_2}} = R = R_L\,}
    \ \Longrightarrow\
    \frac{L_1}{C_2} = R^2\ .
\]

\noindent\textbf{(iii) Despeje.} Tenemos el producto y el cociente de $L_1$ y
$C_2$; multiplicando y dividiendo ambas ecuaciones:
\[
    L_1^2 = (L_1 C_2)\frac{L_1}{C_2} = \frac{4}{\omega_c^2}\,R^2
    \ \Rightarrow\ L_1 = \frac{2R}{\omega_c}\ ,\qquad
    C_2^2 = \frac{L_1 C_2}{L_1/C_2} = \frac{4}{\omega_c^2 R^2}
    \ \Rightarrow\ C_2 = \frac{2}{R\,\omega_c}\ .
\]
Reemplazando $R = R_L = 500\ \Omega$ y $\omega_c = 20.000\ \tfrac{\text{rad}}{\text{s}}$:
\begin{align*}
    L_1 &= \frac{2R}{\omega_c} = \frac{2\cdot 500}{20.000} = 0{,}05\ \text{H} = \boxed{50\ \text{mH}}\\[2pt]
    C_2 &= \frac{2}{R\,\omega_c} = \frac{2}{500\cdot 20.000} = 2\times10^{-7}\ \text{F} = \boxed{200\ \text{nF}}
\end{align*}
La rama serie se reparte en dos brazos de $L_1/2 = 25\ \text{mH}$:

\begin{figure}[H]
    \centering
    % Figura generada con netlist2tikz (n2t render fig_ej02_proto_pb.sch --tikz --no-standalone).
% Netlist fuente (fig_ej02_proto_pb.sch):
%   W in_t 1; right=0.4
%   L1 1 2; right=2, l^=25\,\mathrm{mH}
%   L2 2 3; right=2, l^=25\,\mathrm{mH}
%   W 3 out_t; right=0.4
%   C 2 0; down=1.6, l_=200\,\mathrm{nF}
%   W in_b 0; right=2.4
%   W 0 out_b; right=2.4
%
\begin{tikzpicture}[scale=1.00, transform shape, /tikz/circuitikz/bipoles/length=1.50cm, american currents, american voltages, font={\fontsize{7.5}{9}\selectfont}, voltage dir=RP]
  \coordinate (in_t) at (0,1.92);
  \coordinate (1) at (0.48,1.92);
  \coordinate (2) at (2.88,1.92);
  \coordinate (3) at (5.28,1.92);
  \coordinate (out_t) at (5.76,1.92);
  \coordinate (0) at (2.88,0);
  \coordinate (in_b) at (0,0);
  \coordinate (out_b) at (5.76,0);
  \draw (in_t) to (1);
  \draw (1) to [L, l^=$25\,\mathrm{mH}$, n=L1] (2);
  \draw (2) to [L, l^=$25\,\mathrm{mH}$, n=L2] (3);
  \draw (3) to (out_t);
  \draw (2) to [C, l_=$200\,\mathrm{nF}$, n=C] (0);
  \draw (in_b) to (0);
  \draw (0) to (out_b);
  \draw (in_t) node[ocirc] {};
  \draw (out_t) node[ocirc] {};
  \draw (in_b) node[ocirc] {};
  \draw (out_b) node[ocirc] {};
\end{tikzpicture}

\end{figure}

% --------------------------------------------------------------
\ejercicio{Mismo prototipo, partiendo del circuito normalizado.}

\begin{enunciadobox}
Calcular el ejercicio anterior partiendo del circuito normalizado.
\end{enunciadobox}

El prototipo pasabajos \emph{normalizado} ($\omega_c = 1$, $R_L = 1$) tiene
$\bar L_1 = 2\,$H, $\bar C_2 = 2\,$F (ramas serie $\bar L_1/2 = 1\,$H):

\begin{figure}[H]
    \centering
    % Figura generada con netlist2tikz (n2t render fig_ej03_proto_norm.sch --tikz --no-standalone).
% Netlist fuente (fig_ej03_proto_norm.sch):
%   W in_t 1; right=0.4
%   L1 1 2; right=2, l^=1
%   L2 2 3; right=2, l^=1
%   W 3 out_t; right=0.4
%   C 2 0; down=1.6, l_=2 
%   W in_b 0; right=2.4
%   W 0 out_b; right=2.4
%
\begin{tikzpicture}[scale=1.00, transform shape, /tikz/circuitikz/bipoles/length=1.50cm, american currents, american voltages, font={\fontsize{7.5}{9}\selectfont}, voltage dir=RP]
  \coordinate (in_t) at (0,1.92);
  \coordinate (1) at (0.48,1.92);
  \coordinate (2) at (2.88,1.92);
  \coordinate (3) at (5.28,1.92);
  \coordinate (out_t) at (5.76,1.92);
  \coordinate (0) at (2.88,0);
  \coordinate (in_b) at (0,0);
  \coordinate (out_b) at (5.76,0);
  \draw (in_t) to (1);
  \draw (1) to [L, l^=1, n=L1] (2);
  \draw (2) to [L, l^=1, n=L2] (3);
  \draw (3) to (out_t);
  \draw (2) to [C, l_=2, n=C] (0);
  \draw (in_b) to (0);
  \draw (0) to (out_b);
  \draw (in_t) node[ocirc] {};
  \draw (out_t) node[ocirc] {};
  \draw (in_b) node[ocirc] {};
  \draw (out_b) node[ocirc] {};
\end{tikzpicture}

\end{figure}

\noindent Tomamos las normas (números adimensionales) $a = R_L/1\,\Omega = 500$ y
$N = \omega_c/1\,\tfrac{\text{rad}}{\text{s}} = 20.000$ y desnormalizamos
($\bar L_1,\bar C_2$ ya están en H y F):
\begin{align*}
    L_1 &= \frac{\bar L_1\,a}{N} = \frac{2\cdot 500}{20.000} = 0{,}05\ \text{H} = 50\ \text{mH}
    \quad\Rightarrow\quad \frac{L_1}{2} = 25\ \text{mH}\\[2pt]
    C_2 &= \frac{\bar C_2}{a\,N} = \frac{2}{500\cdot 20.000} = 200\ \text{nF}
\end{align*}
Idéntico al Ejercicio 2, como debía ser.

% --------------------------------------------------------------
\ejercicio{Sección $m$-derivada del prototipo anterior, $\omega_\infty = 22.500\ \tfrac{\text{rad}}{\text{s}}$.}

\begin{enunciadobox}
Tomando el prototipo del ejercicio Nº 2, diseñar la sección $m$-derivada para una
frecuencia de atenuación infinita $\omega_\infty = 22.500\ \tfrac{\text{rad}}{\text{s}}$.
\end{enunciadobox}

Coeficiente de derivación. La relación entre $m$ y la frecuencia de atenuación
infinita es
\[
    m = \sqrt{1-\left(\frac{\omega_c}{\omega_\infty}\right)^2}\ .
\]
Reemplazando $\omega_c = 20.000$ y $\omega_\infty = 22.500\ \tfrac{\text{rad}}{\text{s}}$:
\[
    \left(\frac{\omega_c}{\omega_\infty}\right)^2 = \left(\frac{20.000}{22.500}\right)^2 = 0{,}790
    \ \Rightarrow\
    m = \sqrt{1-0{,}790} = \sqrt{0{,}210} = 0{,}458 \quad (1-m^2 = 0{,}790).
\]
Elementos de la sección $m$-derivada serie (a partir de $L_1 = 50$ mH, $C_2 = 200$ nF):
\begin{align*}
    \text{Rama serie (cada brazo):}\quad & m\,\frac{L_1}{2} = 0{,}458\cdot 25\ \text{mH} = \boxed{11{,}5\ \text{mH}}\\[2pt]
    \text{Derivación (inductor):}\quad & L_1\frac{1-m^2}{4m} = 50\ \text{mH}\cdot\frac{0{,}790}{4\cdot 0{,}458}
    = 50\cdot 0{,}431\ \text{mH} = \boxed{21{,}6\ \text{mH}}\\[2pt]
    \text{Derivación (capacitor):}\quad & m\,C_2 = 0{,}458\cdot 200\ \text{nF} = \boxed{91{,}6\ \text{nF}}
\end{align*}
\noindent Por ser sección completa (T), la rama serie se reparte en \emph{dos}
brazos de $11{,}5$ mH ($z_1/2$ cada uno); la guía la tabula sin dividir, $m L_1 = 23$ mH.
 
\begin{figure}[H]
    \centering
    % Figura generada con netlist2tikz (n2t render fig_ej04_mderiv_pb.sch --tikz --no-standalone).
% Netlist fuente (fig_ej04_mderiv_pb.sch):
%   W in_t 1; right=0.4
%   L1 1 2; right=2, l^=11{,}5\,\mathrm{mH}
%   L2 2 3; right=2, l^=11{,}5\,\mathrm{mH}
%   W 3 out_t; right=0.4
%   L3 2 4; down=1.1, l_=21{,}6\,\mathrm{mH}
%   C 4 0; down=1.1, l_=91{,}6\,\mathrm{nF}
%   W in_b 0; right=2.4
%   W 0 out_b; right=2.4
%
\begin{tikzpicture}[scale=1.00, transform shape, /tikz/circuitikz/bipoles/length=1.50cm, american currents, american voltages, font={\fontsize{7.5}{9}\selectfont}, voltage dir=RP]
  \coordinate (in_t) at (0,2.64);
  \coordinate (1) at (0.48,2.64);
  \coordinate (2) at (2.88,2.64);
  \coordinate (3) at (5.28,2.64);
  \coordinate (out_t) at (5.76,2.64);
  \coordinate (4) at (2.88,1.32);
  \coordinate (0) at (2.88,0);
  \coordinate (in_b) at (0,0);
  \coordinate (out_b) at (5.76,0);
  \draw (in_t) to (1);
  \draw (1) to [L, l^=$11{,}5\,\mathrm{mH}$, n=L1] (2);
  \draw (2) to [L, l^=$11{,}5\,\mathrm{mH}$, n=L2] (3);
  \draw (3) to (out_t);
  \draw (2) to [L, l_=$21{,}6\,\mathrm{mH}$, n=L3] (4);
  \draw (4) to [C, l_=$91{,}6\,\mathrm{nF}$, n=C] (0);
  \draw (in_b) to (0);
  \draw (0) to (out_b);
  \draw (in_t) node[ocirc] {};
  \draw (out_t) node[ocirc] {};
  \draw (in_b) node[ocirc] {};
  \draw (out_b) node[ocirc] {};
\end{tikzpicture}

\end{figure}

\noindent A la frecuencia $\omega_\infty$ la rama derivación entra en resonancia
(cortocircuito a masa) y la salida se anula: aparece el cero de transmisión.

% --------------------------------------------------------------
\ejercicio{Hemisección de adaptación correspondiente ($m = 0{,}6$).}

\begin{enunciadobox}
Calcular la hemi-sección de adaptación correspondiente al ejercicio anterior.
\end{enunciadobox}

La hemisección de adaptación es media sección $L$ con $m = 0{,}6$ (el valor que
hace $z_{0\pi}'$ prácticamente constante en la banda). Se obtiene cortando la
sección completa por la mitad: la rama serie queda como un \emph{único} brazo
($z_1'/2 = m\,z_1/2$), no repartida en dos como en la sección completa, mientras
que la rama derivación \emph{duplica} la impedancia $z_2'$ de la sección completa (su inductor
se duplica y su capacitor se reduce a la mitad). Con $L_1 = 50$ mH, $C_2 = 200$ nF
y $1-m^2 = 0{,}64$:
\begin{align*}
    \text{Rama serie (único brazo):}\quad & m\,\frac{L_1}{2} = 0{,}6\cdot 25\ \text{mH} = \boxed{15\ \text{mH}}\\[2pt]
    \text{Derivación (inductor):}\quad & L_1\frac{1-m^2}{2m} = 50\ \text{mH}\cdot\frac{0{,}64}{1{,}2}
    = \boxed{26{,}7\ \text{mH}}\\[2pt]
    \text{Derivación (capacitor):}\quad & C_2\frac{m}{2} = 200\ \text{nF}\cdot 0{,}3 = \boxed{60\ \text{nF}}
\end{align*}
\noindent La guía tabula la rama serie como $L_1 = 30$ mH (valor sin dividir; el
brazo colocado es la mitad, $15$ mH, según la convención indicada al inicio).

\begin{figure}[H]
    \centering
    % Figura generada con netlist2tikz (n2t render fig_ej05_hemi_pb.sch --tikz --no-standalone).
% Netlist fuente (fig_ej05_hemi_pb.sch):
%   W in_t 1; right=0.4
%   L1 1 c; right=2, l^=15\,\mathrm{mH}
%   W c out_t; right=0.5
%   L2 c 4; down=1.1, l_=26{,}7\,\mathrm{mH}
%   C 4 0; down=1.1, l_=60\,\mathrm{nF}
%   W in_b 0; right=2.4
%   W 0 out_b; right=0.5
%
\begin{tikzpicture}[scale=1.00, transform shape, /tikz/circuitikz/bipoles/length=1.50cm, american currents, american voltages, font={\fontsize{7.5}{9}\selectfont}, voltage dir=RP]
  \coordinate (in_t) at (0,2.64);
  \coordinate (1) at (0.48,2.64);
  \coordinate (c) at (2.88,2.64);
  \coordinate (out_t) at (3.48,2.64);
  \coordinate (4) at (2.88,1.32);
  \coordinate (0) at (2.88,0);
  \coordinate (in_b) at (0,0);
  \coordinate (out_b) at (3.48,0);
  \draw (in_t) to (1);
  \draw (1) to [L, l^=$15\,\mathrm{mH}$, n=L1] (c);
  \draw (c) to (out_t);
  \draw (c) to [L, l_=$26{,}7\,\mathrm{mH}$, n=L2] (4);
  \draw (4) to [C, l_=$60\,\mathrm{nF}$, n=C] (0);
  \draw (in_b) to (0);
  \draw (0) to (out_b);
  \draw (in_t) node[ocirc] {};
  \draw (out_t) node[ocirc] {};
  \draw (in_b) node[ocirc] {};
  \draw (out_b) node[ocirc] {};
\end{tikzpicture}

\end{figure}

\noindent (Las dos hemisecciones $m = 0{,}6$ de las terminaciones equivalen a una
sección $m$-derivada completa, que aporta un cero de transmisión lejano.)

% --------------------------------------------------------------
\ejercicio{Transformaciones de frecuencia a partir del pasabajos.}

\begin{enunciadobox}
A partir del filtro pasabajos obtener, mediante la transformación en frecuencia:
\begin{enumerate}
    \item[a)] un pasaaltos, especificando la transformación de elementos, la
    frecuencia de corte, el ancho de banda y la impedancia característica;
    \item[b)] un pasabanda.
\end{enumerate}
\end{enunciadobox}

\noindent\textbf{a) Pasaaltos --- inversión $s = 1/s'$.}\par
La transformación de los elementos surge de aplicar $s=1/s'$ a las impedancias:
\[
    s'L \;\longrightarrow\; \frac{L}{s} = \frac{1}{s\,(1/L)}\quad(\text{capacitor }1/L)\ ,
    \qquad
    \frac{1}{s'C}\;\longrightarrow\; \frac{s}{C} = s\,(1/C)\quad(\text{inductor }1/C).
\]
Cada \emph{inductor} se vuelve un \emph{capacitor} y cada \emph{capacitor} un
\emph{inductor}. Sobre el prototipo pasabajos (ramas serie $L_1/2$, derivación $C_2$):
\[
    \frac{L_1}{2}\ \longrightarrow\ C_1 = \frac{2}{L_1}\ \text{(cada brazo serie)}\ ,
    \qquad
    C_2\ \longrightarrow\ L_2 = \frac{1}{C_2}\ \text{(derivación)} .
\]

\begin{figure}[H]
    \centering
    % Figura generada con netlist2tikz (n2t render fig_ej06a_proto_pa.sch --tikz --no-standalone).
% Netlist fuente (fig_ej06a_proto_pa.sch):
%   W in_t 1; right=0.4
%   C1 1 2; right=2, l^=\frac{2}{L_1}
%   C2 2 3; right=2, l^=\frac{2}{L_1}
%   W 3 out_t; right=0.4
%   L 2 0; down=1.6, l_=\frac{1}{C_2}
%   W in_b 0; right=2.4
%   W 0 out_b; right=2.4
%
\begin{tikzpicture}[scale=1.00, transform shape, /tikz/circuitikz/bipoles/length=1.50cm, /tikz/circuitikz/label distance=5pt, american currents, american voltages, font={\fontsize{10}{12}\selectfont}, voltage dir=RP]
  \coordinate (in_t) at (0,1.92);
  \coordinate (1) at (0.48,1.92);
  \coordinate (2) at (2.88,1.92);
  \coordinate (3) at (5.28,1.92);
  \coordinate (out_t) at (5.76,1.92);
  \coordinate (0) at (2.88,0);
  \coordinate (in_b) at (0,0);
  \coordinate (out_b) at (5.76,0);
  \draw (in_t) to (1);
  \draw (1) to [C, l^=$\frac{2}{L_1}$, n=C1] (2);
  \draw (2) to [C, l^=$\frac{2}{L_1}$, n=C2] (3);
  \draw (3) to (out_t);
  \draw (2) to [L, l_=$\frac{1}{C_2}$, n=L] (0);
  \draw (in_b) to (0);
  \draw (0) to (out_b);
  \draw (in_t) node[ocirc] {};
  \draw (out_t) node[ocirc] {};
  \draw (in_b) node[ocirc] {};
  \draw (out_b) node[ocirc] {};
\end{tikzpicture}

\end{figure}

\noindent La inversión \emph{preserva la banda}: la característica pasante del
pasabajos ($0<\omega'<\omega_c$) se mapea en la del pasaaltos ($\omega_c<\omega<\infty$).
Llamamos $C_1 = 1/L_1$ y $L_2 = 1/C_2$ al inductor de derivación. Las reactancias
de la rama serie y de la rama derivación son
\[
    z_1 = \frac{1}{j\omega C_1}\ ,\quad X_1 = -\frac{1}{\omega C_1}\ ;\qquad\qquad
    z_2 = j\omega L_2\ ,\quad X_2 = \omega L_2 .
\]
A partir de ellas \emph{deducimos} los tres parámetros pedidos.

\noindent\textbf{Frecuencia de corte.} De la condición de corte $X_1 = -4X_2$:
\[
    \frac{1}{\omega_c C_1} = 4\,\omega_c L_2
    \ \Longrightarrow\ \omega_c^2 = \frac{1}{4\,L_2 C_1}
    \ \Longrightarrow\ \boxed{\;\omega_c = \frac{1}{2\sqrt{L_2\,C_1}}\;}
\]

\noindent\textbf{Impedancia característica.} De la sección T,
\[
    z_{0T} = \sqrt{z_1 z_2 + \frac{z_1^2}{4}}
           = \sqrt{\frac{L_2}{C_1} - \frac{1}{4\,\omega^2 C_1^2}}\ ,
\]
real en la banda pasante. Cuando $\omega\to\infty$ los capacitores serie son
cortocircuitos y el segundo término se anula, quedando la impedancia nominal
\[
    \boxed{\;z_{0T}(\infty) = \sqrt{\frac{L_2}{C_1}} = R\;}\qquad
    \text{(y } z_{0T}=0 \text{ en } \omega=\omega_c\text{, lo que reobtiene el corte).}
\]

\noindent\textbf{Ancho de banda.} El pasaaltos pasa desde $\omega_c$ hasta infinito,
de modo que su banda pasante ---y por tanto su ancho de banda $\text{AB}$--- es
\[
    \boxed{\;\text{AB}: (\omega_c\,;\,\infty)\;}\qquad\text{(teóricamente infinito).}
\]
El Ejercicio 7 aplica estas mismas deducciones a un caso numérico.

\medskip
\noindent\textbf{b) Pasabanda --- transformación cuadrática $s' = s + \dfrac{\omega_0^2}{s}$.}\par
Con $\omega_0 = \sqrt{\omega_1\omega_2}$ (centro geométrico) y ancho de banda
normalizado $=1$. La transformación de los elementos es:
\[
    s'L\ \longrightarrow\ sL + \frac{1}{s\,\frac{1}{\omega_0^2 L}}
    \quad(\text{serie } L\text{--}C),
    \qquad
    s'C\ \longrightarrow\ sC + \frac{1}{s\,\frac{1}{\omega_0^2 C}}
    \quad(\text{paralelo } L\text{--}C).
\]
El inductor serie pasa a un resonante \emph{serie} y el capacitor derivación a un
resonante \emph{paralelo}, ambos sintonizados en $\omega_0$. El prototipo
pasabanda T resulta:

\begin{figure}[H]
    \centering
    % Figura generada con netlist2tikz (n2t render fig_ej06b_proto_pbanda.sch --tikz --no-standalone).
% Netlist fuente (fig_ej06b_proto_pbanda.sch):
%   W in_t 1; right=0.4
%   L1 1 1a; right=1.6, l^=\frac{L_1}{2}
%   C1 1a m1; right=1.6, l^=\frac{2}{\omega_0^2 L_1}
%   W m1 c; right=0.4
%   W c m2; right=0.4
%   C1b m2 2c; right=1.6, l^=\frac{2}{\omega_0^2 L_1}
%   L1b 2c 3; right=1.6, l=\frac{L_1}{2}
%   W 3 out_t; right=0.4
%   W c sp; down=0.6
%   L2 sp 5; down=1.8, offset=-0.5, l_=\frac{1}{\omega_0^2 C_2}
%   C2 sp 5; down=1.8, offset=0.5, l^=C_2
%   W 5 0; down=0.5
%   W in_b 0; right=4.0
%   W 0 out_b; right=4.0
%
\begin{tikzpicture}[scale=1.00, transform shape, /tikz/circuitikz/bipoles/length=1.50cm, american currents, american voltages, font={\fontsize{7.5}{9}\selectfont}, voltage dir=RP]
  \coordinate (in_t) at (0,3.48);
  \coordinate (1) at (0.48,3.48);
  \coordinate (1a) at (2.4,3.48);
  \coordinate (m1) at (4.32,3.48);
  \coordinate (c) at (4.8,3.48);
  \coordinate (m2) at (5.28,3.48);
  \coordinate (2c) at (7.2,3.48);
  \coordinate (3) at (9.12,3.48);
  \coordinate (out_t) at (9.6,3.48);
  \coordinate (sp) at (4.8,2.76);
  \coordinate (spoff1) at (4.2,2.76);
  \coordinate (5) at (4.8,0.6);
  \coordinate (5off2) at (4.2,0.6);
  \coordinate (spoff3) at (5.4,2.76);
  \coordinate (5off4) at (5.4,0.6);
  \coordinate (0) at (4.8,0);
  \coordinate (in_b) at (0,0);
  \coordinate (out_b) at (9.6,0);
  \draw (in_t) to (1);
  \draw (1) to [L, l^=$\frac{L_1}{2}$, n=L1] (1a);
  \draw (1a) to [C, l^=$\frac{2}{\omega_0^2 L_1}$, n=C1] (m1);
  \draw (m1) to (c);
  \draw (c) to (m2);
  \draw (m2) to [C, l^=$\frac{2}{\omega_0^2 L_1}$, n=C1b] (2c);
  \draw (2c) to [L, l_=$\frac{L_1}{2}$, n=L1b] (3);
  \draw (3) to (out_t);
  \draw (c) to (sp);
  \draw (sp) to (spoff1);
  \draw (5) to (5off2);
  \draw (sp) to [open, n=Oanon1] (5);
  \draw (spoff1) to [L, l_=$\frac{1}{\omega_0^2 C_2}$, n=L2] (5off2);
  \draw (sp) to (spoff3);
  \draw (5) to (5off4);
  \draw (sp) to [open, n=Oanon2] (5);
  \draw (spoff3) to [C, l^=$C_2$, n=C2] (5off4);
  \draw (5) to (0);
  \draw (in_b) to (0);
  \draw (0) to (out_b);
  \draw (in_t) node[ocirc] {};
  \draw (out_t) node[ocirc] {};
  \draw (in_b) node[ocirc] {};
  \draw (out_b) node[ocirc] {};
\end{tikzpicture}

\end{figure}

\noindent A partir de la transformación para régimen sinusoidal permanente,
$\bar\omega' = \bar\omega - \dfrac{\bar\omega_0^2}{\bar\omega}$, deducimos ---igual
que en el pasaaltos--- las frecuencias de corte, el ancho de banda y la impedancia
característica del pasabanda.

\noindent\textbf{Frecuencias de corte.} Los cortes del prototipo $\bar\omega' = \pm1$
se mapean resolviendo $\bar\omega - \dfrac{\bar\omega_0^2}{\bar\omega} = \pm1$, o sea
$\bar\omega^2 \mp \bar\omega - \bar\omega_0^2 = 0$:
\[
    \boxed{\;\bar\omega_2 = 0{,}5 + \sqrt{0{,}25 + \bar\omega_0^2}\;}\qquad
    \boxed{\;\bar\omega_1 = -0{,}5 + \sqrt{0{,}25 + \bar\omega_0^2}\;}
\]
(raíces positivas), que verifican $\bar\omega_1\bar\omega_2 = \bar\omega_0^2$
---centro geométrico, $\omega_0 = \sqrt{\omega_1\omega_2}$--- y $\bar\omega_2-\bar\omega_1 = 1$.

\noindent\textbf{Ancho de banda.} Es la diferencia de cortes; normalizado vale $1$
y real coincide con la norma de frecuencia:
\[
    \boxed{\;\text{AB} = \omega_2-\omega_1 = N\;}
\]

\noindent\textbf{Impedancia característica.} Conviene ver \emph{por qué} alcanza con
sustituir $\bar\omega'\to\bar\omega-\bar\omega_0^2/\bar\omega$ en la $z_{0T}$ del
prototipo. La impedancia característica depende de la frecuencia \emph{sólo a través de
las reactancias} de las dos ramas, pues $z_{0T}=\sqrt{z_1 z_2+\frac{z_1^2}{4}}$. Basta
entonces ver cómo se transforma cada reactancia.

\noindent\underline{Rama serie.} En el prototipo es un inductor, $z_1=j\bar\omega'\bar L$.
Transformada, es $\bar L$ en serie con $\bar C=1/(\bar\omega_0^2\bar L)$; su impedancia a
la frecuencia $\bar\omega$ del pasabanda es
\[
    z_1(\bar\omega) = j\bar\omega\bar L + \frac{1}{j\bar\omega\bar C}
    = j\bar\omega\bar L + \frac{\bar\omega_0^2\bar L}{j\bar\omega}
    = j\bar L\left(\bar\omega-\frac{\bar\omega_0^2}{\bar\omega}\right)
    = j\bar L\,\bar\omega' .
\]
Es decir: la rama serie transformada tiene en $\bar\omega$ \emph{la misma reactancia} que
la del prototipo en la frecuencia imagen $\bar\omega'$.

\noindent\underline{Rama derivación.} En el prototipo es un capacitor; trabajamos con su
admitancia, $y_2=j\bar\omega'\bar C_2$. Transformada, es $\bar C_2$ en paralelo con
$\bar L=1/(\bar\omega_0^2\bar C_2)$, de admitancia
\[
    y_2(\bar\omega) = j\bar\omega\bar C_2 + \frac{1}{j\bar\omega\bar L}
    = j\bar C_2\left(\bar\omega-\frac{\bar\omega_0^2}{\bar\omega}\right)
    = j\bar C_2\,\bar\omega' ,
\]
o sea $z_2(\bar\omega)=z_2(\bar\omega')\big|_{\text{proto}}$ también.

\noindent\underline{Conclusión.} La $z_{0T}$ del \emph{prototipo pasabajos} la 
calculamos en el Ejercicio~2:
\[
    z_{0T}^{\,\text{proto}}(\omega') = \sqrt{\frac{L_1}{C_2}}\,\sqrt{1-\left(\frac{\omega'}{\omega_c}\right)^{\!2}}
    \;\xrightarrow[\;\omega_c\to1\;]{\text{normalizando}}\;
    R\sqrt{1-\bar\omega'^{\,2}}\qquad\Big(R=\sqrt{\tfrac{L_1}{C_2}}\Big).
\]
Como las dos reactancias del pasabanda en $\bar\omega$ coinciden con las del prototipo en
$\bar\omega'=\bar\omega-\bar\omega_0^2/\bar\omega$, \emph{cualquier} función de ellas
---en particular esa $z_{0T}$--- hereda la sustitución $\bar\omega'\to\bar\omega-\bar\omega_0^2/\bar\omega$:
\[
    z_{0T}^{\,\text{PB}}(\bar\omega) = z_{0T}^{\,\text{proto}}(\bar\omega')
    \;\Longrightarrow\;
    \boxed{\;z_{0T}(\omega)=R\sqrt{1-\left(\bar\omega-\frac{\bar\omega_0^2}{\bar\omega}\right)^{\!2}}\;}
\]
con $R=\sqrt{L_1/C_2}=R_L$. Esto \emph{es} lo que significa una ``transformación de
frecuencia'': la curva reactancia--frecuencia del prototipo ---y con ella $z_{0T}$---
se vuelve a leer a través del cambio $\bar\omega\mapsto\bar\omega'$. Es real en la
banda: vale la nominal $R$ en el centro $\omega_0$ ($\bar\omega'=0$) y se anula en los
cortes $\omega_1,\omega_2$ ($\bar\omega'=\pm1$). El producto $z_1 z_2=k^2=L_1/C_2$ se
conserva, por eso la impedancia nominal sigue siendo $R_L$.

% --------------------------------------------------------------
\ejercicio{Prototipo pasaaltos dado: $C = 10^{-7}\,$F, hallar $L$ y $\omega_c$.}

\begin{enunciadobox}
El circuito de la figura representa el prototipo de un filtro pasa altos imagen,
donde $C = 10^{-7}\ \text{F}$:
\begin{figure}[H]
    \centering
    % Figura generada con netlist2tikz (n2t render fig_ej07_enunciado.sch --tikz --no-standalone).
% Netlist fuente (fig_ej07_enunciado.sch):
%   W in_t 1; right=0.4
%   C1 1 2; right=2, l^=C
%   C2 2 3; right=2, l^=C
%   W 3 out_t; right=0.4
%   L 2 0; down=1.6, l_=L
%   W in_b 0; right=2.4
%   W 0 out_b; right=2.4
%
\begin{tikzpicture}[scale=1.00, transform shape, /tikz/circuitikz/bipoles/length=1.50cm, american currents, american voltages, font={\fontsize{7.5}{9}\selectfont}, voltage dir=RP]
  \coordinate (in_t) at (0,1.92);
  \coordinate (1) at (0.48,1.92);
  \coordinate (2) at (2.88,1.92);
  \coordinate (3) at (5.28,1.92);
  \coordinate (out_t) at (5.76,1.92);
  \coordinate (0) at (2.88,0);
  \coordinate (in_b) at (0,0);
  \coordinate (out_b) at (5.76,0);
  \draw (in_t) to (1);
  \draw (1) to [C, l^=C, n=C1] (2);
  \draw (2) to [C, l^=C, n=C2] (3);
  \draw (3) to (out_t);
  \draw (2) to [L, l_=L, n=L] (0);
  \draw (in_b) to (0);
  \draw (0) to (out_b);
  \draw (in_t) node[ocirc] {};
  \draw (out_t) node[ocirc] {};
  \draw (in_b) node[ocirc] {};
  \draw (out_b) node[ocirc] {};
\end{tikzpicture}

\end{figure}
Calcular:
\begin{enumerate}
    \item[a)] el valor del elemento ``$L$'' sabiendo que la impedancia
    característica $z_{0T}(\infty) = 100\ \Omega$ cuando $\omega\to\infty$;
    \item[b)] la frecuencia de corte $\omega_c$.
\end{enumerate}
\end{enunciadobox}

La rama serie tiene un capacitor $C$ en cada brazo de la T (la impedancia de cada brazo es
$z_1/2 = \dfrac{1}{j\omega C}$), de modo que la impedancia  $z_1 = \dfrac{2}{j\omega C}$, y la derivación es $z_2 = j\omega L$.
La impedancia característica de la sección T:
\[
    z_{0T} = \sqrt{z_1 z_2 + \frac{z_1^2}{4}}
           = \sqrt{\frac{2L}{C} + \frac{1}{4}\left(\frac{2}{j\omega C}\right)^{\!2}}
           = \sqrt{\frac{2L}{C} - \frac{1}{\omega^2 C^2}}\ .
\]

\noindent\textbf{a)} Cuando $\omega\to\infty$ el segundo término se anula:
\[
    z_{0T}(\infty) = \sqrt{\frac{2L}{C}} = 100\ \Omega
    \ \Rightarrow\ \frac{2L}{C} = 10^4
    \ \Rightarrow\ L = \frac{10^4\,C}{2} = 5000\cdot 10^{-7}\ \text{H}
    = \boxed{500\ \mu\text{H}}\ .
\]

\noindent\textbf{b)} La frecuencia de corte sale de la condición $X_1 = -4X_2$, con
las reactancias de la sección pasaaltos $X_1 = -\dfrac{2}{\omega C}$ (rama serie) y
$X_2 = \omega L$ (rama derivación):
\[
    -\frac{2}{\omega_c C} = -4\,\omega_c L
    \ \Rightarrow\ \omega_c^2 = \frac{1}{2LC}
    \ \Rightarrow\ \omega_c = \frac{1}{\sqrt{2LC}}
    = \frac{1}{\sqrt{2\cdot 5\times10^{-4}\cdot 10^{-7}}}
    = \frac{1}{\sqrt{10^{-10}}} = \boxed{100.000\ \tfrac{\text{rad}}{\text{s}}}\ .
\]

\begin{figure}[H]
    \centering
    % Figura generada con netlist2tikz (n2t render fig_ej07_sol.sch --tikz --no-standalone).
% Netlist fuente (fig_ej07_sol.sch):
%   W in_t 1; right=0.4
%   C1 1 2; right=2, l^=0{,}1\,\mu\mathrm{F}
%   C2 2 3; right=2, l^=0{,}1\,\mu\mathrm{F}
%   W 3 out_t; right=0.4
%   L 2 0; down=1.6, l_=500\,\mu\mathrm{H}
%   W in_b 0; right=2.4
%   W 0 out_b; right=2.4
%
\begin{tikzpicture}[scale=1.00, transform shape, /tikz/circuitikz/bipoles/length=1.50cm, american currents, american voltages, font={\fontsize{7.5}{9}\selectfont}, voltage dir=RP]
  \coordinate (in_t) at (0,1.92);
  \coordinate (1) at (0.48,1.92);
  \coordinate (2) at (2.88,1.92);
  \coordinate (3) at (5.28,1.92);
  \coordinate (out_t) at (5.76,1.92);
  \coordinate (0) at (2.88,0);
  \coordinate (in_b) at (0,0);
  \coordinate (out_b) at (5.76,0);
  \draw (in_t) to (1);
  \draw (1) to [C, l^=$0{,}1\,\mu\mathrm{F}$, n=C1] (2);
  \draw (2) to [C, l^=$0{,}1\,\mu\mathrm{F}$, n=C2] (3);
  \draw (3) to (out_t);
  \draw (2) to [L, l_=$500\,\mu\mathrm{H}$, n=L] (0);
  \draw (in_b) to (0);
  \draw (0) to (out_b);
  \draw (in_t) node[ocirc] {};
  \draw (out_t) node[ocirc] {};
  \draw (in_b) node[ocirc] {};
  \draw (out_b) node[ocirc] {};
\end{tikzpicture}

\end{figure}

% --------------------------------------------------------------
\ejercicio{Hemisección de la sección $m$-derivada pasaaltos normalizada.}

\begin{enunciadobox}
Dado el circuito $m$-derivada normalizado de un filtro pasaalto (designados como
$C_1'$, $L'$, $C_2'$), calcular la hemisección. Determinar los valores normalizados.
\begin{figure}[H]
    \centering
    % Figura generada con netlist2tikz (n2t render fig_ej08_enunciado.sch --tikz --no-standalone).
% Netlist fuente (fig_ej08_enunciado.sch):
%   W in_t 1; right=0.4
%   C1 1 2; right=2, l^=C_1'
%   C2 2 3; right=2, l^=C_1'
%   W 3 out_t; right=0.4
%   L 2 4; down=1.1, l_=L'
%   C3 4 0; down=1.1, l_=C_2'
%   W in_b 0; right=2.4
%   W 0 out_b; right=2.4
%
\begin{tikzpicture}[scale=1.00, transform shape, /tikz/circuitikz/bipoles/length=1.50cm, american currents, american voltages, font={\fontsize{7.5}{9}\selectfont}, voltage dir=RP]
  \coordinate (in_t) at (0,2.64);
  \coordinate (1) at (0.48,2.64);
  \coordinate (2) at (2.88,2.64);
  \coordinate (3) at (5.28,2.64);
  \coordinate (out_t) at (5.76,2.64);
  \coordinate (4) at (2.88,1.32);
  \coordinate (0) at (2.88,0);
  \coordinate (in_b) at (0,0);
  \coordinate (out_b) at (5.76,0);
  \draw (in_t) to (1);
  \draw (1) to [C, l^=$C_1'$, n=C1] (2);
  \draw (2) to [C, l^=$C_1'$, n=C2] (3);
  \draw (3) to (out_t);
  \draw (2) to [L, l_=L', n=L] (4);
  \draw (4) to [C, l_=$C_2'$, n=C3] (0);
  \draw (in_b) to (0);
  \draw (0) to (out_b);
  \draw (in_t) node[ocirc] {};
  \draw (out_t) node[ocirc] {};
  \draw (in_b) node[ocirc] {};
  \draw (out_b) node[ocirc] {};
\end{tikzpicture}

\end{figure}
\[
    C_1' = \frac{2}{m L_1}\ ;\qquad
    C_2' = \frac{4m}{L_1(1-m^2)}\ ;\qquad
    L' = \frac{1}{m C_2}
\]
Tener en cuenta que los valores del pasabajos normalizado son:
$L_1 = 2$; $C_2 = 2$ y $\omega_\infty = \sqrt{1-m^2}$.
\end{enunciadobox}

Partimos del pasabajos normalizado ($L_1 = 2$, $C_2 = 2$) y de $m = 0{,}6$
(valor de adaptación). Los elementos de la sección $m$-derivada pasaaltos son
los dados en el enunciado:
\[
    C_1' = \frac{2}{m L_1} = \frac{2}{0{,}6\cdot 2} = 1{,}67\ \text{F},\quad
    L'   = \frac{1}{m C_2} = \frac{1}{0{,}6\cdot 2} = 0{,}83\ \text{H},\quad
    C_2' = \frac{4m}{L_1(1-m^2)} = \frac{2{,}4}{2\cdot 0{,}64} = 1{,}875\ \text{F}.
\]

\noindent La hemisección se obtiene cortando la sección T
$m$-derivada por su rama derivación. El brazo serie queda igual ($z_1/2 = C_1'$),
y la rama derivación \emph{duplica su impedancia} ($2z_2$): el inductor se
duplica y el capacitor se reduce a la mitad. Por lo tanto:
\[
    \boxed{\,C_s = C_1' = \frac{2}{m L_1} = 1{,}67\ \text{F}\,}\qquad
    \boxed{\,L_d = 2L' = \frac{2}{m C_2} = 1{,}67\ \text{H}\,}\qquad
    \boxed{\,C_d = \frac{C_2'}{2} = \frac{2m}{(1-m^2)L_1} = 0{,}94\ \text{F}\,}
\]

\begin{figure}[H]
    \centering
    % Figura generada con netlist2tikz (n2t render fig_ej08_hemi_pa.sch --tikz --no-standalone).
% Netlist fuente (fig_ej08_hemi_pa.sch):
%   W in_t 1; right=0.4
%   L1 1 4; down=1.1, l_=L_d
%   C1 4 0; down=1.1, l_=C_d
%   Cs 1 c; right=2, l^=C_s
%   W c out_t; right=0.4
%   W in_b 0; right=0.4
%   W 0 out_b; right=2.4
%
\begin{tikzpicture}[scale=1.00, transform shape, /tikz/circuitikz/bipoles/length=1.50cm, american currents, american voltages, font={\fontsize{7.5}{9}\selectfont}, voltage dir=RP]
  \coordinate (in_t) at (0,2.64);
  \coordinate (1) at (0.48,2.64);
  \coordinate (4) at (0.48,1.32);
  \coordinate (0) at (0.48,0);
  \coordinate (c) at (2.88,2.64);
  \coordinate (out_t) at (3.36,2.64);
  \coordinate (in_b) at (0,0);
  \coordinate (out_b) at (3.36,0);
  \draw (in_t) to (1);
  \draw (1) to [L, l_=$L_d$, n=L1] (4);
  \draw (4) to [C, l_=$C_d$, n=C1] (0);
  \draw (1) to [C, l^=$C_s$, n=Cs] (c);
  \draw (c) to (out_t);
  \draw (in_b) to (0);
  \draw (0) to (out_b);
  \draw (in_t) node[ocirc] {};
  \draw (out_t) node[ocirc] {};
  \draw (in_b) node[ocirc] {};
  \draw (out_b) node[ocirc] {};
\end{tikzpicture}

\end{figure}

\noindent Los valores quedan en forma \emph{normalizada} (Farad y Henry);
para un caso real se desnormalizan con las normas $a = R_L/1\,\Omega$ y
$N = \omega_c/1\,\tfrac{\text{rad}}{\text{s}}$.

% --------------------------------------------------------------
\ejercicio{Filtro completo pasabanda: $R_L = 600\,\Omega$, $\omega_1 = 45.000$, $\omega_2 = 90.000$, $\omega_{\infty 1} = 38.000\ \tfrac{\text{rad}}{\text{s}}$.}

\begin{enunciadobox}
Calcular un filtro completo pasabanda con $R_L = 600\ \Omega$;
$\omega_1 = 45.000\ \tfrac{\text{rad}}{\text{s}}$;
$\omega_2 = 90.000\ \tfrac{\text{rad}}{\text{s}}$;
$\omega_{\infty 1} = 38.000\ \tfrac{\text{rad}}{\text{s}}$.
\end{enunciadobox}

\noindent\underline{Normas y constantes.} Para pasabanda la norma de frecuencia
es el ancho de banda y la de impedancia $R_L$, ambas como números adimensionales:
\[
    N = \frac{\omega_2-\omega_1}{1\,\tfrac{\text{rad}}{\text{s}}} = 45.000\ ,\qquad
    a = \frac{R_L}{1\,\Omega} = 600 .
\]
\[
    \omega_0 = \sqrt{\omega_1\omega_2} = \sqrt{45.000\cdot 90.000} = 63.640\ \tfrac{\text{rad}}{\text{s}}\ ,\qquad
    \bar\omega_0^{\,2} = \frac{\omega_1\omega_2}{N^2} = \frac{4{,}05\times10^9}{(45.000)^2} = 2{,}0 .
\]
Factores de desnormalización ($\bar L,\bar C$ en H y F):
$\;L = \bar L\,\dfrac{a}{N} = \bar L\cdot 13{,}33\ \text{mH}\;$ y
$\;C = \dfrac{\bar C}{a\,N} = \bar C\cdot 37{,}04\ \text{nF}$.

\medskip
\noindent\textbf{Prototipo.} Tomamos directamente el modelo normalizado del
prototipo pasabanda~T del paquete de modelos del Cap.~4 (\emph{no} rehacemos la
transformación de frecuencia: ya está hecha allí, de una vez para siempre). En
función de $\bar\omega_0$, el modelo da: rama serie ---cada brazo--- un resonante
\emph{serie} con $\bar L = 1$ y $\bar C = \dfrac{1}{\bar\omega_0^2}$; rama derivación
un resonante \emph{paralelo} con $\bar L = \dfrac{1}{2\bar\omega_0^2}$ y $\bar C = 2$.
Con $\bar\omega_0^2 = 2$:
\[
\begin{array}{lll}
    \text{Rama serie (cada brazo):} & \bar L = 1,\;\; \bar C = \tfrac{1}{\bar\omega_0^2} = 0{,}5 & (\text{resonante serie})\\[3pt]
    \text{Derivación:} & \bar L = \tfrac{1}{2\bar\omega_0^2} = 0{,}25,\;\; \bar C = 2 & (\text{resonante paralelo})
\end{array}
\]
Desnormalizamos con $L = \bar L\,\tfrac{a}{N} = \bar L\cdot 13{,}33$ mH y
$C = \dfrac{\bar C}{a\,N} = \bar C\cdot 37{,}04$ nF:
\[
    \begin{array}{llll}
        L_1 = 1\cdot 13{,}33 = 13\ \text{mH} &\quad& C_1 = 0{,}5\cdot 37{,}04 = 18{,}5\ \text{nF} &\ (\text{rama serie})\\[2pt]
        L_2 = 0{,}25\cdot 13{,}33 = 3{,}3\ \text{mH} && C_2 = 2\cdot 37{,}04 = 74\ \text{nF} &\ (\text{derivación})
    \end{array}
\]

\begin{figure}[H]
    \centering
    % Figura generada con netlist2tikz (n2t render fig_ej09_proto.sch --tikz --no-standalone).
% Netlist fuente (fig_ej09_proto.sch):
%   W in_t 1; right=0.4
%   L1 1 1a; right=1.5, l^=13\,\mathrm{mH}
%   C1 1a m1; right=1.5, l^=18{,}5\,\mathrm{nF}
%   W m1 c; right=0.4
%   W c m2; right=0.4
%   C1b m2 2c; right=1.5, l^=18{,}5\,\mathrm{nF}
%   L1b 2c 3; right=1.5, l=13\,\mathrm{mH}
%   W 3 out_t; right=0.4
%   W c sp; down=0.5
%   L2 sp 5; down=1.6, offset=-0.45, l_=3{,}3\,\mathrm{mH}
%   C2 sp 5; down=1.6, offset=0.45, l^=74\,\mathrm{nF}
%   W 5 0; down=0.5
%   W in_b 0; right=3.8
%   W 0 out_b; right=3.8
%
\begin{tikzpicture}[scale=1.00, transform shape, /tikz/circuitikz/bipoles/length=1.50cm, american currents, american voltages, font={\fontsize{7.5}{9}\selectfont}, voltage dir=RP]
  \coordinate (in_t) at (0,3.12);
  \coordinate (1) at (0.48,3.12);
  \coordinate (1a) at (2.28,3.12);
  \coordinate (m1) at (4.08,3.12);
  \coordinate (c) at (4.56,3.12);
  \coordinate (m2) at (5.04,3.12);
  \coordinate (2c) at (6.84,3.12);
  \coordinate (3) at (8.64,3.12);
  \coordinate (out_t) at (9.12,3.12);
  \coordinate (sp) at (4.56,2.52);
  \coordinate (spoff1) at (4.02,2.52);
  \coordinate (5) at (4.56,0.6);
  \coordinate (5off2) at (4.02,0.6);
  \coordinate (spoff3) at (5.1,2.52);
  \coordinate (5off4) at (5.1,0.6);
  \coordinate (0) at (4.56,0);
  \coordinate (in_b) at (0,0);
  \coordinate (out_b) at (9.12,0);
  \draw (in_t) to (1);
  \draw (1) to [L, l^=$13\,\mathrm{mH}$, n=L1] (1a);
  \draw (1a) to [C, l^=$18{,}5\,\mathrm{nF}$, n=C1] (m1);
  \draw (m1) to (c);
  \draw (c) to (m2);
  \draw (m2) to [C, l^=$18{,}5\,\mathrm{nF}$, n=C1b] (2c);
  \draw (2c) to [L, l_=$13\,\mathrm{mH}$, n=L1b] (3);
  \draw (3) to (out_t);
  \draw (c) to (sp);
  \draw (sp) to (spoff1);
  \draw (5) to (5off2);
  \draw (sp) to [open, n=Oanon1] (5);
  \draw (spoff1) to [L, l_=$3{,}3\,\mathrm{mH}$, n=L2] (5off2);
  \draw (sp) to (spoff3);
  \draw (5) to (5off4);
  \draw (sp) to [open, n=Oanon2] (5);
  \draw (spoff3) to [C, l^=$74\,\mathrm{nF}$, n=C2] (5off4);
  \draw (5) to (0);
  \draw (in_b) to (0);
  \draw (0) to (out_b);
  \draw (in_t) node[ocirc] {};
  \draw (out_t) node[ocirc] {};
  \draw (in_b) node[ocirc] {};
  \draw (out_b) node[ocirc] {};
\end{tikzpicture}

\end{figure}

\noindent\textbf{Sección $m$-derivada.}

\noindent\underline{Coeficiente $m$.} La frecuencia de cero $\omega_{\infty 1} = 38.000$
se normaliza con $N$: $\bar\omega_\infty = 38.000/45.000 = 0{,}844$
($\bar\omega_\infty^2 = 0{,}713$). En el pasabajos de partida el cero cae en
$\bar\omega'_\infty = 1/\sqrt{1-m^2}$, y ese $\bar\omega'_\infty$ es la \emph{imagen}
por la transformación cuadrática de la frecuencia real $\bar\omega_\infty$, esto es
$\bar\omega'_\infty = \bar\omega_\infty - \dfrac{\bar\omega_0^2}{\bar\omega_\infty}$.
Como la característica del pasabajos es función \emph{par} de $\bar\omega'$ (el cero
existe en $\pm\bar\omega'_\infty$), la igualdad vale \emph{en módulo}; acá hace falta:
$\omega_{\infty 1}$ está por debajo de la banda y la imagen da negativa
($0{,}844 - 2/0{,}844 = -1{,}525$). Igualamos en módulo y elevamos al cuadrado:
\[
    \frac{1}{\sqrt{1-m^2}} = \left|\bar\omega_\infty - \frac{\bar\omega_0^2}{\bar\omega_\infty}\right|
    = \frac{\left|\bar\omega_\infty^2-\bar\omega_0^2\right|}{\bar\omega_\infty}
    \;\Rightarrow\;
    1-m^2 = \frac{\bar\omega_\infty^2}{\left(\bar\omega_\infty^2-\bar\omega_0^2\right)^2},
\]
que es exactamente la fórmula del Cap.~4 (Modelos Normalizados):
\[
    m = \sqrt{1-\frac{\bar\omega_\infty^2}{\left(\bar\omega_\infty^2-\bar\omega_0^2\right)^2}}
      = \sqrt{1-\frac{0{,}713}{(0{,}713-2)^2}}
      = \sqrt{1-\frac{0{,}713}{1{,}656}} = \sqrt{0{,}569} = 0{,}755 \quad (1-m^2 = 0{,}431).
\]

\noindent\underline{Elementos.} Tomamos el modelo normalizado del $m$-derivado
pasabanda~T (Cap.~4). En función de $m$ y $\bar\omega_0$: rama serie ---cada brazo---
un resonante \emph{serie} con $\bar L=m$ y $\bar C=\dfrac{1}{\bar\omega_0^2 m}$; en la
derivación, un bloque resonante \emph{paralelo} ($\bar L=\dfrac{1}{2\bar\omega_0^2 m}$,
$\bar C=2m$) en serie con un bloque resonante \emph{serie}
($\bar C=\dfrac{2m}{\bar\omega_0^2(1-m^2)}$, $\bar L=\dfrac{1-m^2}{2m}$). Con
$m=0{,}755$ y $\bar\omega_0^2=2$:
\[
    \begin{array}{lll}
        \bar L=m=0{,}755 & \bar C=\tfrac{1}{\bar\omega_0^2 m}=0{,}662 & (\text{rama serie})\\[3pt]
        \bar L=\tfrac{1}{2\bar\omega_0^2 m}=0{,}331 & \bar C=2m=1{,}51 & (\text{bloque paralelo})\\[3pt]
        \bar C=\tfrac{2m}{\bar\omega_0^2(1-m^2)}=1{,}752 & \bar L=\tfrac{1-m^2}{2m}=0{,}285 & (\text{bloque serie})
    \end{array}
\]
Desnormalizando con $L=\bar L\,\tfrac{a}{N}=\bar L\cdot 13{,}33$ mH y $C=\dfrac{\bar C}{a\,N}=\bar C\cdot 37{,}04$ nF:
\[
    \begin{array}{lll}
        L_1 = 0{,}755\cdot 13{,}33 = 10\ \text{mH} & C_1 = 0{,}662\cdot 37{,}04 = 24{,}5\ \text{nF} & (\text{rama serie})\\[3pt]
        L_2 = 0{,}331\cdot 13{,}33 = 4{,}4\ \text{mH} & C_2 = 1{,}51\cdot 37{,}04 = 56\ \text{nF} & (\text{bloque paralelo})\\[3pt]
        L_3 = 0{,}285\cdot 13{,}33 = 3{,}8\ \text{mH} & C_3 = 1{,}752\cdot 37{,}04 = 65\ \text{nF} & (\text{bloque serie})
    \end{array}
\]

\begin{figure}[H]
    \centering
    % Figura generada con netlist2tikz (n2t render fig_ej09_mderiv.sch --tikz --no-standalone).
% Netlist fuente (fig_ej09_mderiv.sch):
%   W in_t 1; right=0.4
%   L1 1 1a; right=1.5, l^=10\,\mathrm{mH}
%   C1 1a m1; right=1.5, l^=24{,}5\,\mathrm{nF}
%   W m1 c; right=0.4
%   W c m2; right=0.4
%   C1b m2 2c; right=1.5, l^=24{,}5\,\mathrm{nF}
%   L1b 2c 3; right=1.5, l=10\,\mathrm{mH}
%   W 3 out_t; right=0.4
%   W c sp; down=0.5
%   L2 sp t1; down=1.5, offset=-0.45, l_=4{,}4\,\mathrm{mH}
%   C2 sp t1; down=1.5, offset=0.45, l^=56\,\mathrm{nF}
%   L3 t1 t2; down=1.0, l_=3{,}8\,\mathrm{mH}
%   C3 t2 0; down=1.0, l_=65\,\mathrm{nF}
%   W in_b 0; right=3.8
%   W 0 out_b; right=3.8
%
\begin{tikzpicture}[scale=1.00, transform shape, /tikz/circuitikz/bipoles/length=1.50cm, american currents, american voltages, font={\fontsize{7.5}{9}\selectfont}, voltage dir=RP]
  \coordinate (in_t) at (0,4.8);
  \coordinate (1) at (0.48,4.8);
  \coordinate (1a) at (2.28,4.8);
  \coordinate (m1) at (4.08,4.8);
  \coordinate (c) at (4.56,4.8);
  \coordinate (m2) at (5.04,4.8);
  \coordinate (2c) at (6.84,4.8);
  \coordinate (3) at (8.64,4.8);
  \coordinate (out_t) at (9.12,4.8);
  \coordinate (sp) at (4.56,4.2);
  \coordinate (spoff1) at (4.02,4.2);
  \coordinate (t1) at (4.56,2.4);
  \coordinate (t1off2) at (4.02,2.4);
  \coordinate (spoff3) at (5.1,4.2);
  \coordinate (t1off4) at (5.1,2.4);
  \coordinate (t2) at (4.56,1.2);
  \coordinate (0) at (4.56,0);
  \coordinate (in_b) at (0,0);
  \coordinate (out_b) at (9.12,0);
  \draw (in_t) to (1);
  \draw (1) to [L, l^=$10\,\mathrm{mH}$, n=L1] (1a);
  \draw (1a) to [C, l^=$24{,}5\,\mathrm{nF}$, n=C1] (m1);
  \draw (m1) to (c);
  \draw (c) to (m2);
  \draw (m2) to [C, l^=$24{,}5\,\mathrm{nF}$, n=C1b] (2c);
  \draw (2c) to [L, l_=$10\,\mathrm{mH}$, n=L1b] (3);
  \draw (3) to (out_t);
  \draw (c) to (sp);
  \draw (sp) to (spoff1);
  \draw (t1) to (t1off2);
  \draw (sp) to [open, n=Oanon1] (t1);
  \draw (spoff1) to [L, l_=$4{,}4\,\mathrm{mH}$, n=L2] (t1off2);
  \draw (sp) to (spoff3);
  \draw (t1) to (t1off4);
  \draw (sp) to [open, n=Oanon2] (t1);
  \draw (spoff3) to [C, l^=$56\,\mathrm{nF}$, n=C2] (t1off4);
  \draw (t1) to [L, l_=$3{,}8\,\mathrm{mH}$, n=L3] (t2);
  \draw (t2) to [C, l_=$65\,\mathrm{nF}$, n=C3] (0);
  \draw (in_b) to (0);
  \draw (0) to (out_b);
  \draw (in_t) node[ocirc] {};
  \draw (out_t) node[ocirc] {};
  \draw (in_b) node[ocirc] {};
  \draw (out_b) node[ocirc] {};
\end{tikzpicture}

\end{figure}

\noindent\textbf{Hemisección de adaptación} ($m = 0{,}6$). Las dos hemisecciones de
adaptación equivalen a una sección $m$-derivada completa de $m=0{,}6$ (Cap.~4); por
eso la obtenemos \emph{bisecando} el modelo $m$-derivado pasabanda con $m=0{,}6$: la
rama serie queda en un \emph{único} brazo (el del modelo) y la rama derivación
\emph{duplica} su impedancia (en cada bloque resonante el inductor se duplica y el
capacitor se reduce a la mitad). Con $\bar\omega_0^2=2$:
\[
    \begin{array}{lll}
        \bar L=m=0{,}6 & \bar C=\tfrac{1}{\bar\omega_0^2 m}=0{,}833 & (\text{rama serie, un brazo})\\[3pt]
        \bar L=2\cdot\tfrac{1}{2\bar\omega_0^2 m}=0{,}833 & \bar C=\tfrac{2m}{2}=0{,}6 & (\text{deriv.\ paralelo, }2z)\\[3pt]
        \bar L=2\cdot\tfrac{1-m^2}{2m}=1{,}067 & \bar C=\tfrac{1}{2}\tfrac{2m}{\bar\omega_0^2(1-m^2)}=0{,}469 & (\text{deriv.\ serie, }2z)
    \end{array}
\]
Desnormalizando con $L=\bar L\,\tfrac{a}{N}=\bar L\cdot 13{,}33$ mH y $C=\dfrac{\bar C}{a\,N}=\bar C\cdot 37{,}04$ nF:
\[
    \begin{array}{lll}
        L_1 = 0{,}6\cdot 13{,}33 = 8\ \text{mH} & C_1 = 0{,}833\cdot 37{,}04 = 31\ \text{nF} & (\text{rama serie})\\[3pt]
        L_2 = 0{,}833\cdot 13{,}33 = 11{,}1\ \text{mH} & C_2 = 0{,}6\cdot 37{,}04 = 22{,}2\ \text{nF} & (\text{paralelo})\\[3pt]
        L_3 = 1{,}067\cdot 13{,}33 = 14{,}2\ \text{mH} & C_3 = 0{,}469\cdot 37{,}04 = 17{,}4\ \text{nF} & (\text{serie})
    \end{array}
\]

\begin{figure}[H]
    \centering
    % Figura generada con netlist2tikz (n2t render fig_ej09_hemi.sch --tikz --no-standalone).
% Netlist fuente (fig_ej09_hemi.sch):
%   W in_t 1; right=0.4
%   L1 1 1a; right=1.5, l^=8\,\mathrm{mH}
%   C1 1a c; right=1.5, l^=31\,\mathrm{nF}
%   W c out_t; right=0.6
%   W c sp; down=0.5
%   L2 sp t1; down=1.5, offset=-0.45, l_=11{,}1\,\mathrm{mH}
%   C2 sp t1; down=1.5, offset=0.45, l^=22{,}2\,\mathrm{nF}
%   L3 t1 t2; down=1.0, l_=14{,}2\,\mathrm{mH}
%   C3 t2 0; down=1.0, l_=17{,}4\,\mathrm{nF}
%   W in_b 0; right=3.4
%   W 0 out_b; right=0.6
%
\begin{tikzpicture}[scale=1.00, transform shape, /tikz/circuitikz/bipoles/length=1.50cm, american currents, american voltages, font={\fontsize{7.5}{9}\selectfont}, voltage dir=RP]
  \coordinate (in_t) at (0,4.8);
  \coordinate (1) at (0.48,4.8);
  \coordinate (1a) at (2.28,4.8);
  \coordinate (c) at (4.08,4.8);
  \coordinate (out_t) at (4.8,4.8);
  \coordinate (sp) at (4.08,4.2);
  \coordinate (spoff1) at (3.54,4.2);
  \coordinate (t1) at (4.08,2.4);
  \coordinate (t1off2) at (3.54,2.4);
  \coordinate (spoff3) at (4.62,4.2);
  \coordinate (t1off4) at (4.62,2.4);
  \coordinate (t2) at (4.08,1.2);
  \coordinate (0) at (4.08,0);
  \coordinate (in_b) at (0,0);
  \coordinate (out_b) at (4.8,0);
  \draw (in_t) to (1);
  \draw (1) to [L, l^=$8\,\mathrm{mH}$, n=L1] (1a);
  \draw (1a) to [C, l^=$31\,\mathrm{nF}$, n=C1] (c);
  \draw (c) to (out_t);
  \draw (c) to (sp);
  \draw (sp) to (spoff1);
  \draw (t1) to (t1off2);
  \draw (sp) to [open, n=Oanon1] (t1);
  \draw (spoff1) to [L, l_=$11{,}1\,\mathrm{mH}$, n=L2] (t1off2);
  \draw (sp) to (spoff3);
  \draw (t1) to (t1off4);
  \draw (sp) to [open, n=Oanon2] (t1);
  \draw (spoff3) to [C, l^=$22{,}2\,\mathrm{nF}$, n=C2] (t1off4);
  \draw (t1) to [L, l_=$14{,}2\,\mathrm{mH}$, n=L3] (t2);
  \draw (t2) to [C, l_=$17{,}4\,\mathrm{nF}$, n=C3] (0);
  \draw (in_b) to (0);
  \draw (0) to (out_b);
  \draw (in_t) node[ocirc] {};
  \draw (out_t) node[ocirc] {};
  \draw (in_b) node[ocirc] {};
  \draw (out_b) node[ocirc] {};
\end{tikzpicture}

\end{figure}

\begin{nota}
\textit{Verificación:} en cada par $L$-$C$ producto de la transformación debe
cumplirse la resonancia en $\omega_0$; p.\,ej. $L_2 C_2 = 11{,}1\,\text{mH}\cdot 22{,}2\,\text{nF}$
da $\omega = 1/\sqrt{L_2 C_2} \approx 63.700 \approx \omega_0$, y lo mismo para
$L_3 C_3 = 14{,}2\,\text{mH}\cdot 17{,}4\,\text{nF}$.
\end{nota}

\noindent\textbf{Filtro completo:} se conecta en cascada con adaptación imagen
\;\textsf{[Hemisec. $m{=}0{,}6$] -- [Prototipo] -- [$m$-derivada] -- [Hemisec. $m{=}0{,}6$]}\;
entre $R_G$ y $R_L$, según el esquema del Ejercicio 1.

% --------------------------------------------------------------
\ejercicio{Filtro completo suprimebanda: $R_L = 75\,\Omega$, $f_1 = 1.000$, $f_2 = 3.500$, $f_{\infty 2} = 3.100\ \text{Hz}$.}

\begin{enunciadobox}
Calcular un filtro completo suprimebanda con $R_L = 75\ \Omega$;
$f_1 = 1.000\ \text{Hz}$; $f_2 = 3.500\ \text{Hz}$; $f_{\infty 2} = 3.100\ \text{Hz}$.
\end{enunciadobox}

\noindent\underline{Datos en pulsación.}
$\omega_1 = 2\pi f_1 = 6.283$, $\omega_2 = 2\pi f_2 = 21.991$,
$\omega_{\infty 2} = 2\pi f_{\infty 2} = 19.478\ \tfrac{\text{rad}}{\text{s}}$.
\[
    N = \frac{\omega_2-\omega_1}{1\,\tfrac{\text{rad}}{\text{s}}} = 2\pi(3.500-1.000) = 15.708\ ,\qquad
    a = \frac{R_L}{1\,\Omega} = 75 ,
\]
\[
    \omega_0 = \sqrt{\omega_1\omega_2} = 11.756\ \tfrac{\text{rad}}{\text{s}}\ ,\qquad
    \bar\omega_0^{\,2} = \frac{\omega_1\omega_2}{N^2} = 0{,}560 .
\]
Factores de desnormalización ($\bar L,\bar C$ en H y F):
$\;L = \bar L\,\dfrac{a}{N} = \bar L\cdot 4{,}775\ \text{mH}\;$ y
$\;C = \dfrac{\bar C}{a\,N} = \bar C\cdot 848{,}8\ \text{nF}$.

\medskip
\noindent\textbf{Prototipo.} Tomamos el modelo normalizado del prototipo
suprimebanda~T del Cap.~4 (sin rehacer la transformación). En función de
$\bar\omega_0$: rama serie ---cada brazo--- un resonante \emph{paralelo} con
$\bar C=1$ y $\bar L=\dfrac{1}{\bar\omega_0^2}$; rama derivación un resonante
\emph{serie} con $\bar L=0{,}5$ y $\bar C=\dfrac{2}{\bar\omega_0^2}$. Con
$\bar\omega_0^2=0{,}560$:
\[
\begin{array}{lll}
    \text{Rama serie (cada brazo):} & \bar C=1,\;\; \bar L=\tfrac{1}{\bar\omega_0^2}=1{,}786 & (\text{resonante paralelo})\\[3pt]
    \text{Derivación:} & \bar L=0{,}5,\;\; \bar C=\tfrac{2}{\bar\omega_0^2}=3{,}571 & (\text{resonante serie})
\end{array}
\]
Desnormalizando con $L=\bar L\,\tfrac{a}{N}=\bar L\cdot 4{,}775$ mH y $C=\dfrac{\bar C}{a\,N}=\bar C\cdot 848{,}8$ nF:
\[
    \begin{array}{llll}
        L_1 = 1{,}786\cdot 4{,}775 = 8{,}5\ \text{mH} &\quad& C_1 = 1\cdot 848{,}8 = 849\ \text{nF} &\ (\text{serie: } L_1\,\|\,C_1)\\[2pt]
        L_2 = 0{,}5\cdot 4{,}775 = 2{,}4\ \text{mH} && C_2 = 3{,}571\cdot 848{,}8 = 3.030\ \text{nF} &\ (\text{derivación serie})
    \end{array}
\]

\begin{figure}[H]
    \centering
    % Figura generada con netlist2tikz (n2t render fig_ej10_proto.sch --tikz --no-standalone).
% Netlist fuente (fig_ej10_proto.sch):
%   W in_t 1; right=0.5
%   L1 1 m1; right=2, offset=0.35, l^=8{,}5\,\mathrm{mH}
%   C1 1 m1; right=2, offset=-0.35, l_=849\,\mathrm{nF}
%   W m1 c; right=0.45
%   W c m2; right=0.45
%   L2 m2 3; right=2, offset=0.35, l^=8{,}5\,\mathrm{mH}
%   C2 m2 3; right=2, offset=-0.35, l_=849\,\mathrm{nF}
%   W 3 out_t; right=0.5
%   L3 c 4; down=1.2, l_=2{,}4\,\mathrm{mH}
%   C3 4 0; down=1.2, l_=3.030\,\mathrm{nF}
%   W in_b 0; right=3.4
%   W 0 out_b; right=3.4
%
\begin{tikzpicture}[scale=1.00, transform shape, /tikz/circuitikz/bipoles/length=1.50cm, american currents, american voltages, font={\fontsize{7.5}{9}\selectfont}, voltage dir=RP]
  \coordinate (in_t) at (0.54,2.88);
  \coordinate (1) at (1.14,2.88);
  \coordinate (1off1) at (1.14,3.3);
  \coordinate (m1) at (3.54,2.88);
  \coordinate (m1off2) at (3.54,3.3);
  \coordinate (1off3) at (1.14,2.46);
  \coordinate (m1off4) at (3.54,2.46);
  \coordinate (c) at (4.08,2.88);
  \coordinate (m2) at (4.62,2.88);
  \coordinate (m2off5) at (4.62,3.3);
  \coordinate (3) at (7.02,2.88);
  \coordinate (3off6) at (7.02,3.3);
  \coordinate (m2off7) at (4.62,2.46);
  \coordinate (3off8) at (7.02,2.46);
  \coordinate (out_t) at (7.62,2.88);
  \coordinate (4) at (4.08,1.44);
  \coordinate (0) at (4.08,0);
  \coordinate (in_b) at (0,0);
  \coordinate (out_b) at (8.16,0);
  \draw (in_t) to (1);
  \draw (1) to (1off1);
  \draw (m1) to (m1off2);
  \draw (1) to [open, n=Oanon1] (m1);
  \draw (1off1) to [L, l^=$8{,}5\,\mathrm{mH}$, n=L1] (m1off2);
  \draw (1) to (1off3);
  \draw (m1) to (m1off4);
  \draw (1) to [open, n=Oanon2] (m1);
  \draw (1off3) to [C, l_=$849\,\mathrm{nF}$, n=C1] (m1off4);
  \draw (m1) to (c);
  \draw (c) to (m2);
  \draw (m2) to (m2off5);
  \draw (3) to (3off6);
  \draw (m2) to [open, n=Oanon3] (3);
  \draw (m2off5) to [L, l^=$8{,}5\,\mathrm{mH}$, n=L2] (3off6);
  \draw (m2) to (m2off7);
  \draw (3) to (3off8);
  \draw (m2) to [open, n=Oanon4] (3);
  \draw (m2off7) to [C, l_=$849\,\mathrm{nF}$, n=C2] (3off8);
  \draw (3) to (out_t);
  \draw (c) to [L, l_=$2{,}4\,\mathrm{mH}$, n=L3] (4);
  \draw (4) to [C, l_=$3.030\,\mathrm{nF}$, n=C3] (0);
  \draw (in_b) to (0);
  \draw (0) to (out_b);
  \draw (in_t) node[ocirc] {};
  \draw (out_t) node[ocirc] {};
  \draw (in_b) node[ocirc] {};
  \draw (out_b) node[ocirc] {};
\end{tikzpicture}

\end{figure}

\noindent\textbf{Sección $m$-derivada.} Con $\omega_{\infty 2}$ en la banda suprimida,
$\bar\omega_\infty = 19.478/15.708 = 1{,}240$. Como el suprimebanda viene del
\emph{pasaaltos}, la relación de partida es $\bar\omega'_\infty = \sqrt{1-m^2}$, que con la
transformación cuadrática $\bar\omega'_\infty = \bar\omega_\infty - \bar\omega_0^2/\bar\omega_\infty$ da:
\[
\begin{aligned}
    \sqrt{1-m^2} = \bar\omega_\infty - \frac{\bar\omega_0^2}{\bar\omega_\infty}
    \ &\Rightarrow\
    m = \sqrt{1-\left(\bar\omega_\infty - \frac{\bar\omega_0^2}{\bar\omega_\infty}\right)^2}\\
    &= \sqrt{1-\left(1{,}240 - \frac{0{,}560}{1{,}240}\right)^2}
      = \sqrt{1-(0{,}788)^2} = \sqrt{0{,}379} = 0{,}615 \quad (1-m^2 = 0{,}621).
\end{aligned}
\]
Tomamos el modelo normalizado del $m$-derivado suprimebanda~T (Cap.~4). En función
de $m$ y $\bar\omega_0$: rama serie ---cada brazo--- un resonante \emph{paralelo} con
$\bar C=\dfrac{1}{m}$ y $\bar L=\dfrac{m}{\bar\omega_0^2}$; en la derivación, un bloque
resonante \emph{serie} ($\bar L=\dfrac{1}{2m}$, $\bar C=\dfrac{2m}{\bar\omega_0^2}$) en
serie con un bloque resonante \emph{paralelo} ($\bar C=\dfrac{2m}{1-m^2}$,
$\bar L=\dfrac{1-m^2}{2m\bar\omega_0^2}$). Con $m=0{,}615$ y $\bar\omega_0^2=0{,}560$:
\[
    \begin{array}{lll}
        \bar C=\tfrac{1}{m}=1{,}626 & \bar L=\tfrac{m}{\bar\omega_0^2}=1{,}098 & (\text{rama serie})\\[3pt]
        \bar L=\tfrac{1}{2m}=0{,}813 & \bar C=\tfrac{2m}{\bar\omega_0^2}=2{,}196 & (\text{bloque serie})\\[3pt]
        \bar C=\tfrac{2m}{1-m^2}=1{,}981 & \bar L=\tfrac{1-m^2}{2m\bar\omega_0^2}=0{,}901 & (\text{bloque paralelo})
    \end{array}
\]
Desnormalizando con $L=\bar L\,\tfrac{a}{N}=\bar L\cdot 4{,}775$ mH y $C=\dfrac{\bar C}{a\,N}=\bar C\cdot 848{,}8$ nF:
\[
    \begin{array}{lll}
        L_1 = 1{,}098\cdot 4{,}775 = 5{,}3\ \text{mH} & C_1 = 1{,}626\cdot 848{,}8 = 1.380\ \text{nF} & (\text{rama serie})\\[3pt]
        L_2 = 0{,}813\cdot 4{,}775 = 3{,}9\ \text{mH} & C_2 = 2{,}196\cdot 848{,}8 = 1.860\ \text{nF} & (\text{bloque serie})\\[3pt]
        L_3 = 0{,}901\cdot 4{,}775 = 4{,}3\ \text{mH} & C_3 = 1{,}981\cdot 848{,}8 = 1.680\ \text{nF} & (\text{bloque paralelo})
    \end{array}
\]

\begin{figure}[H]
    \centering
    % Figura generada con netlist2tikz (n2t render fig_ej10_mderiv.sch --tikz --no-standalone).
% Netlist fuente (fig_ej10_mderiv.sch):
%   W in_t 1; right=0.5
%   L1 1 m1; right=2, offset=0.35, l^=5{,}3\,\mathrm{mH}
%   C1 1 m1; right=2, offset=-0.35, l_=1.380\,\mathrm{nF}
%   W m1 c; right=0.45
%   W c m2; right=0.45
%   L2 m2 3; right=2, offset=0.35, l^=5{,}3\,\mathrm{mH}
%   C2 m2 3; right=2, offset=-0.35, l_=1.380\,\mathrm{nF}
%   W 3 out_t; right=0.5
%   L3 c 4; down=1.1, l_=3{,}9\,\mathrm{mH}
%   C3 4 5; down=1.1, l_=1.860\,\mathrm{nF}
%   L4 5 6; down=1.6, offset=-0.55, l_=4{,}3\,\mathrm{mH}
%   C4 5 6; down=1.6, offset=0.55, l^=1.680\,\mathrm{nF}
%   W 6 0; down=0.5
%   W in_b 0; right=3.4
%   W 0 out_b; right=3.4
%
\begin{tikzpicture}[scale=1.00, transform shape, /tikz/circuitikz/bipoles/length=1.50cm, american currents, american voltages, font={\fontsize{7.5}{9}\selectfont}, voltage dir=RP]
  \coordinate (in_t) at (0.54,5.16);
  \coordinate (1) at (1.14,5.16);
  \coordinate (1off1) at (1.14,5.58);
  \coordinate (m1) at (3.54,5.16);
  \coordinate (m1off2) at (3.54,5.58);
  \coordinate (1off3) at (1.14,4.74);
  \coordinate (m1off4) at (3.54,4.74);
  \coordinate (c) at (4.08,5.16);
  \coordinate (m2) at (4.62,5.16);
  \coordinate (m2off5) at (4.62,5.58);
  \coordinate (3) at (7.02,5.16);
  \coordinate (3off6) at (7.02,5.58);
  \coordinate (m2off7) at (4.62,4.74);
  \coordinate (3off8) at (7.02,4.74);
  \coordinate (out_t) at (7.62,5.16);
  \coordinate (4) at (4.08,3.84);
  \coordinate (5) at (4.08,2.52);
  \coordinate (5off9) at (3.42,2.52);
  \coordinate (6) at (4.08,0.6);
  \coordinate (6off10) at (3.42,0.6);
  \coordinate (5off11) at (4.74,2.52);
  \coordinate (6off12) at (4.74,0.6);
  \coordinate (0) at (4.08,0);
  \coordinate (in_b) at (0,0);
  \coordinate (out_b) at (8.16,0);
  \draw (in_t) to (1);
  \draw (1) to (1off1);
  \draw (m1) to (m1off2);
  \draw (1) to [open, n=Oanon1] (m1);
  \draw (1off1) to [L, l^=$5{,}3\,\mathrm{mH}$, n=L1] (m1off2);
  \draw (1) to (1off3);
  \draw (m1) to (m1off4);
  \draw (1) to [open, n=Oanon2] (m1);
  \draw (1off3) to [C, l_=$1.380\,\mathrm{nF}$, n=C1] (m1off4);
  \draw (m1) to (c);
  \draw (c) to (m2);
  \draw (m2) to (m2off5);
  \draw (3) to (3off6);
  \draw (m2) to [open, n=Oanon3] (3);
  \draw (m2off5) to [L, l^=$5{,}3\,\mathrm{mH}$, n=L2] (3off6);
  \draw (m2) to (m2off7);
  \draw (3) to (3off8);
  \draw (m2) to [open, n=Oanon4] (3);
  \draw (m2off7) to [C, l_=$1.380\,\mathrm{nF}$, n=C2] (3off8);
  \draw (3) to (out_t);
  \draw (c) to [L, l_=$3{,}9\,\mathrm{mH}$, n=L3] (4);
  \draw (4) to [C, l_=$1.860\,\mathrm{nF}$, n=C3] (5);
  \draw (5) to (5off9);
  \draw (6) to (6off10);
  \draw (5) to [open, n=Oanon5] (6);
  \draw (5off9) to [L, l_=$4{,}3\,\mathrm{mH}$, n=L4] (6off10);
  \draw (5) to (5off11);
  \draw (6) to (6off12);
  \draw (5) to [open, n=Oanon6] (6);
  \draw (5off11) to [C, l^=$1.680\,\mathrm{nF}$, n=C4] (6off12);
  \draw (6) to (0);
  \draw (in_b) to (0);
  \draw (0) to (out_b);
  \draw (in_t) node[ocirc] {};
  \draw (out_t) node[ocirc] {};
  \draw (in_b) node[ocirc] {};
  \draw (out_b) node[ocirc] {};
\end{tikzpicture}

\end{figure}

\noindent\textbf{Hemisección de adaptación} ($m = 0{,}6$). La obtenemos
\emph{bisecando} el modelo $m$-derivado suprimebanda con $m=0{,}6$: la rama serie
queda en un \emph{único} brazo (el del modelo) y la rama derivación \emph{duplica} su
impedancia (en cada bloque, inductor $\times2$ y capacitor $\div2$). Con
$\bar\omega_0^2 = 0{,}560$:
\[
    \begin{array}{lll}
        \bar C=\tfrac{1}{m}=1{,}667 & \bar L=\tfrac{m}{\bar\omega_0^2}=1{,}071 & (\text{rama serie, un brazo})\\[3pt]
        \bar C=\tfrac{1}{2}\tfrac{2m}{1-m^2}=0{,}938 & \bar L=2\tfrac{1-m^2}{2m\bar\omega_0^2}=1{,}905 & (\text{deriv.\ paralelo, }2z)\\[3pt]
        \bar L=2\tfrac{1}{2m}=1{,}667 & \bar C=\tfrac{1}{2}\tfrac{2m}{\bar\omega_0^2}=1{,}071 & (\text{deriv.\ serie, }2z)
    \end{array}
\]
Desnormalizando con $L=\bar L\,\tfrac{a}{N}=\bar L\cdot 4{,}775$ mH y $C=\dfrac{\bar C}{a\,N}=\bar C\cdot 848{,}8$ nF:
\[
    \begin{array}{lll}
        L_1 = 1{,}071\cdot 4{,}775 = 5{,}1\ \text{mH} & C_1 = 1{,}667\cdot 848{,}8 = 1.415\ \text{nF} & (\text{rama serie})\\[3pt]
        L_2 = 1{,}905\cdot 4{,}775 = 9{,}1\ \text{mH} & C_2 = 0{,}938\cdot 848{,}8 = 796\ \text{nF} & (\text{deriv.\ paralelo})\\[3pt]
        L_3 = 1{,}667\cdot 4{,}775 = 8{,}0\ \text{mH} & C_3 = 1{,}071\cdot 848{,}8 = 910\ \text{nF} & (\text{deriv.\ serie})
    \end{array}
\]

\begin{nota}
\textit{Verificación:} cada par $L$-$C$ resuena en
$\omega_0=\sqrt{\omega_1\omega_2}\approx 11.756\ \tfrac{\text{rad}}{\text{s}}$, como debe ser.
\end{nota}

\noindent\textbf{Filtro completo:} se conecta en cascada con adaptación imagen
\;\textsf{[Hemisec. $m{=}0{,}6$] -- [Prototipo] -- [$m$-derivada] -- [Hemisec. $m{=}0{,}6$]}\;
entre $R_G$ y $R_L = 75\ \Omega$, según el esquema del Ejercicio 1.

\vspace{1ex}
\begin{center}\color{utnblue}$\blacklozenge$\end{center}

\end{document}
